% Lesson 59 — Vocabulary: Words and expressions related to difficulties with the use of ICTs at work or for learning — Anglais Tle SH — Cours signé OnBuch+
% Programme officiel MINESEC. Compiler avec : tectonic EN59-vocabulary-words-and-expressions-related-to-difficultie.tex
\newcommand{\DOCMATIERE}{Anglais}
\newcommand{\DOCNIVEAU}{Tle SH}
\newcommand{\DOCMODULE}{Chapter 5 : Media and Communication}
\newcommand{\DOCLECON}{EN59}
\newcommand{\DOCTITRE}{Lesson 59 — Vocabulary: Words and expressions related to difficulties with the use of ICTs at work or for learning}
% OnBuch+ — préambule « cours signés » ENRICHI (compilé avec tectonic / XeLaTeX)
% Système visuel « Pop » d'OnBuch+ : fond crème (jamais blanc), bordures encre
% épaisses, ombres « dures » décalées, titres Archivo Black en capitales.
% Version MATHÉMATIQUES : graphes (pgfplots/tikz), boîtes pédagogiques
% (définition, propriété, méthode, attention, le savais-tu, exercice résolu,
% exercice, TP, bilan), page de garde et signature OnBuch+ ancrée en bas.
%
% Macros attendues dans main.tex AVANT \input{preamble} :
%   \DOCMATIERE  \DOCNIVEAU  \DOCMODULE  \DOCLECON (numéro)  \DOCTITRE
\documentclass[11pt]{article}

\usepackage{fontspec}
\usepackage[english,french]{babel} % français = langue principale ; l'anglais via \eng{..} / textebox
\usepackage[a4paper,margin=2cm,top=3.4cm,bottom=2.8cm,headheight=1.4cm,footskip=1.3cm]{geometry}
\usepackage{amsmath,amssymb,amsfonts}
\usepackage{mathtools} % \xmapsto, \coloneqq... souvent utilisés par les modèles
\usepackage{cancel} % simplifications barrées (bilans de forces)
% \abs{z} (module, valeur absolue) et \norm{u} (norme), utilisés par les modèles
\providecommand{\abs}{}\let\abs\relax\DeclarePairedDelimiter{\abs}{\lvert}{\rvert}
\providecommand{\norm}{}\let\norm\relax\DeclarePairedDelimiter{\norm}{\lVert}{\rVert}
\usepackage[table]{xcolor} % [table] : fournit aussi \rowcolors (alternance de lignes)
\usepackage{pagecolor}
\usepackage{tikz}
\usetikzlibrary{arrows.meta,positioning,calc,shapes.geometric,shapes.misc,shapes.arrows,decorations.pathmorphing,decorations.pathreplacing,decorations.markings,patterns,shadows,mindmap,trees,fit,backgrounds,matrix,intersections,angles,quotes}
\usepackage{pgfplots}
\usepackage[version=4]{mhchem} % équations nucléaires \ce{^{235}_{92}U}
\usepackage[european,siunitx]{circuitikz} % schémas électriques
\pgfplotsset{compat=1.18}
\usepgfplotslibrary{fillbetween,groupplots} % aires entre courbes (calcul intégral)
\usepackage{enumitem}
\usepackage{fancyhdr}
% [most] charge toutes les bibliothèques tcolorbox (listings, minted, raster,
% documentation...) et gonfle inutilement la mémoire TeX sur un document long
% (~300 boîtes) : on ne charge que ce qui est réellement utilisé.
\usepackage[skins,breakable]{tcolorbox}
\usepackage{graphicx}
\usepackage{colortbl}
\usepackage{booktabs}
\usepackage{tabularx}
\usepackage{diagbox} % case d'en-tête diagonale des tableaux à double entrée
% Colonnes extensibles centrée / gauche / droite (C, L, R), souvent utilisées par les modèles
\newcolumntype{C}{>{\centering\arraybackslash}X}
\newcolumntype{L}{>{\raggedright\arraybackslash}X}
\newcolumntype{R}{>{\raggedleft\arraybackslash}X}
\usepackage{array}
\usepackage{multicol}
\usepackage{multirow}
\usepackage{siunitx}
% parse-numbers=false : accepte aussi \SI{2,5 \times 10^{-3}}{...}
\sisetup{locale=FR,output-decimal-marker={,},group-minimum-digits=4,parse-numbers=false}
\DeclareSIUnit\ohm{\ensuremath{\symup{\Omega}}} % U+2126 absent de la police
\DeclareSIUnit\division{div} % réglages d'oscilloscope (ms/div, V/div)
\DeclareSIUnit\tr{tr} % tours (vitesse de rotation, ex. \SI{1500}{\tr\per\minute})
\DeclareSIUnit\molar{M} % concentration molaire (ex. \SI{3}{\molar}, \SI{1}{\micro\molar})
\DeclareSIUnit\an{an} % année (le modèle confond parfois avec \year, un compteur TeX) — ex. \SI{5}{\kilo\meter\per\an}
\DeclareSIUnit\FCFA{FCFA} % franc CFA (ex. \SI{500}{\FCFA\per\kilo\gram})
\DeclareSIUnit\degreeBrix{\textdegree Brix} % teneur en sucre d'un sirop/jus (réfractométrie)
\DeclareSIUnit\month{mois} % le modèle utilise parfois \month, un compteur TeX, comme unité de durée
\usepackage{qrcode}
% \degree : commande fréquemment utilisée par les modèles pour un angle ou une
% température en dehors de \SI{}{} (non fournie par les paquets chargés).
\providecommand{\degree}{^{\circ}}
% \qed (fin de démonstration), utilisé par les modèles sans amsthm
\providecommand{\qed}{\hfill$\square$}
% \barre : double barre des valeurs interdites dans un tableau de variations (tkz-tab)
\providecommand{\barre}{\ensuremath{\|}}
% environnement proof (amsthm non chargé) utilisé par les modèles
\ifdefined\proof\else\newenvironment{proof}[1][Démonstration]{\par\noindent\textit{#1.}\ }{\qed\par}\fi
\usepackage{lastpage}
\usepackage{hyperref}
\hypersetup{hidelinks}

% ============================================================
% Palette « Pop » OnBuch+ — mêmes hex que la classe P (plus_ui.dart)
% ============================================================
\definecolor{poporange}{HTML}{F59321}
\definecolor{popdark}{HTML}{E07A0C}
\definecolor{popcream}{HTML}{FFF6E8}
\definecolor{popink}{HTML}{1C1714}
\definecolor{poppurple}{HTML}{7A5AE0}
\definecolor{popgreen}{HTML}{2E9E6A}
\definecolor{poppink}{HTML}{E0457B}
\definecolor{popblue}{HTML}{2F7DE1}
\definecolor{popgold}{HTML}{C98A12}
\definecolor{popmuted}{HTML}{8A7F76}
\definecolor{popline}{HTML}{EADFD2}
\definecolor{popgray}{HTML}{8A7F76}
\colorlet{popgrey}{popgray}
% Alias tolérants (le modèle nomme parfois les couleurs différemment)
\colorlet{popwhite}{white}
\colorlet{popblack}{popink}
\colorlet{popred}{poppink}
\colorlet{popyellow}{popgold}
% Teintes claires (fonds de tableaux/encadrés) — le modèle nomme parfois
% les couleurs avec un suffixe L (ex. popblueL) sans les avoir définies.
\colorlet{poporangeL}{poporange!20}
\colorlet{popdarkL}{popdark!20}
\colorlet{poppurpleL}{poppurple!20}
\colorlet{popgreenL}{popgreen!20}
\colorlet{poppinkL}{poppink!20}
\colorlet{popblueL}{popblue!20}
\colorlet{popgoldL}{popgold!20}
\colorlet{popredL}{popred!20}
\colorlet{popgrayL}{popgray!20}
% Teintes claires pour fonds de tableaux / zones
\colorlet{popblueL}{popblue!10!white}
\colorlet{poporangeL}{poporange!14!white}
\colorlet{popgreenL}{popgreen!12!white}
\colorlet{poppurpleL}{poppurple!12!white}
\colorlet{poppinkL}{poppink!12!white}
\colorlet{popgoldL}{popgold!14!white}

\pagecolor{popcream}

% ============================================================
% Typographie
% ============================================================
\defaultfontfeatures{Ligatures=TeX}
\setmainfont{PlusJakartaSans-Medium.ttf}[Path=../../../fonts/,BoldFont=PlusJakartaSans-Bold.ttf]
% Maths en sans-serif (Fira Math) avec chiffres et lettres droites de Plus
% Jakarta, pour que les formules se fondent dans le texte.
\usepackage[math-style=ISO,bold-style=ISO]{unicode-math}
\setmathfont{FiraMath-Regular.otf}[Path=../../../fonts/]
\setmathfont{PlusJakartaSans-Medium.ttf}[Path=../../../fonts/,range={up/num,up/{Latin,latin},\mathup}]
\usepackage{icomma} % virgule décimale sans espace en mode math
% Caractères mathématiques Unicode absents des polices de texte (Plus Jakarta,
% Archivo Black) : rendus par leur équivalent mathématique, en texte comme en
% formule — sinon ils s'affichent en carré vide (ex. « Arithmétique dans ℤ »).
\usepackage{newunicodechar}
\newunicodechar{ℝ}{\ensuremath{\mathbb{R}}}
\newunicodechar{ℤ}{\ensuremath{\mathbb{Z}}}
\newunicodechar{ℂ}{\ensuremath{\mathbb{C}}}
\newunicodechar{ℕ}{\ensuremath{\mathbb{N}}}
\newunicodechar{ℚ}{\ensuremath{\mathbb{Q}}}
\newunicodechar{𝔻}{\ensuremath{\mathbb{D}}}
\newunicodechar{α}{\ensuremath{\alpha}}
\newunicodechar{β}{\ensuremath{\beta}}
\newunicodechar{γ}{\ensuremath{\gamma}}
\newunicodechar{δ}{\ensuremath{\delta}}
\newunicodechar{ε}{\ensuremath{\varepsilon}}
\newunicodechar{θ}{\ensuremath{\theta}}
\newunicodechar{λ}{\ensuremath{\lambda}}
\newunicodechar{μ}{\ensuremath{\mu}}
\newunicodechar{σ}{\ensuremath{\sigma}}
\newunicodechar{τ}{\ensuremath{\tau}}
\newunicodechar{φ}{\ensuremath{\varphi}}
\newunicodechar{ψ}{\ensuremath{\psi}}
\newunicodechar{ω}{\ensuremath{\omega}}
\newunicodechar{Σ}{\ensuremath{\Sigma}}
% majuscules produites par \MakeUppercase dans les titres (α-aminés -> Α)
\newunicodechar{Α}{\ensuremath{\alpha}}
\newunicodechar{Β}{\ensuremath{\beta}}
\newunicodechar{Γ}{\ensuremath{\Gamma}}
\newunicodechar{Δ}{\ensuremath{\Delta}}
\newunicodechar{Ε}{\ensuremath{\varepsilon}}
\newunicodechar{Θ}{\ensuremath{\Theta}}
\newunicodechar{Λ}{\ensuremath{\Lambda}}
\newunicodechar{Μ}{\ensuremath{\mu}}
\newunicodechar{Τ}{\ensuremath{\tau}}
\newunicodechar{Φ}{\ensuremath{\Phi}}
\newunicodechar{Ψ}{\ensuremath{\Psi}}
\newunicodechar{Ω}{\ensuremath{\Omega}}
\newunicodechar{↦}{\ensuremath{\mapsto}}
\newunicodechar{∈}{\ensuremath{\in}}
\newunicodechar{∉}{\ensuremath{\notin}}
\newunicodechar{∧}{\ensuremath{\wedge}}
\newunicodechar{⇔}{\ensuremath{\Leftrightarrow}}
\newunicodechar{⟺}{\ensuremath{\Longleftrightarrow}}
\newunicodechar{⇒}{\ensuremath{\Rightarrow}}
\newunicodechar{⟹}{\ensuremath{\Longrightarrow}}
\newunicodechar{∘}{\ensuremath{\circ}}
\newunicodechar{∪}{\ensuremath{\cup}}
\newunicodechar{∩}{\ensuremath{\cap}}
\newunicodechar{≡}{\ensuremath{\equiv}}
\newunicodechar{⊂}{\ensuremath{\subset}}
\newunicodechar{⊥}{\ensuremath{\perp}}
\newunicodechar{∃}{\ensuremath{\exists}}
\newunicodechar{∀}{\ensuremath{\forall}}
\newunicodechar{∛}{\ensuremath{\sqrt[3]{\,}}}
\newunicodechar{∜}{\ensuremath{\sqrt[4]{\,}}}
\newunicodechar{⃗}{\ensuremath{{}^{\to}}}
\newunicodechar{ⁿ}{\ifmmode^{n}\else\textsuperscript{\ensuremath{n}}\fi}
\newunicodechar{ˣ}{\ifmmode^{x}\else\textsuperscript{\ensuremath{x}}\fi}
\newunicodechar{ᵇ}{\ifmmode^{b}\else\textsuperscript{\ensuremath{b}}\fi}
\newunicodechar{ʳ}{\ifmmode^{r}\else\textsuperscript{\ensuremath{r}}\fi}
\newunicodechar{ᵘ}{\ifmmode^{u}\else\textsuperscript{\ensuremath{u}}\fi}
\newunicodechar{ʸ}{\ifmmode^{y}\else\textsuperscript{\ensuremath{y}}\fi}
\newunicodechar{ᵃ}{\ifmmode^{a}\else\textsuperscript{\ensuremath{a}}\fi}
\newunicodechar{ᵉ}{\ifmmode^{e}\else\textsuperscript{\ensuremath{e}}\fi}
\newunicodechar{ᵏ}{\ifmmode^{k}\else\textsuperscript{\ensuremath{k}}\fi}
\newunicodechar{ᵖ}{\ifmmode^{p}\else\textsuperscript{\ensuremath{p}}\fi}
\newunicodechar{ᴺ}{\ifmmode^{N}\else\textsuperscript{\ensuremath{N}}\fi}
\newunicodechar{ᶜ}{\ifmmode^{c}\else\textsuperscript{\ensuremath{c}}\fi}
\newunicodechar{ʰ}{\ifmmode^{h}\else\textsuperscript{\ensuremath{h}}\fi}
\newunicodechar{ᵠ}{\ifmmode^{q}\else\textsuperscript{\ensuremath{q}}\fi}
\newunicodechar{ₙ}{\ifmmode_{n}\else\textsubscript{\ensuremath{n}}\fi}
\newunicodechar{ₐ}{\ifmmode_{a}\else\textsubscript{\ensuremath{a}}\fi}
\newunicodechar{ᵢ}{\ifmmode_{i}\else\textsubscript{\ensuremath{i}}\fi}
\newunicodechar{ᵣ}{\ifmmode_{r}\else\textsubscript{\ensuremath{r}}\fi}
\newunicodechar{ₖ}{\ifmmode_{k}\else\textsubscript{\ensuremath{k}}\fi}
\newunicodechar{ᵦ}{\ifmmode_{\beta}\else\textsubscript{\ensuremath{\beta}}\fi}
\newunicodechar{ₓ}{\ifmmode_{x}\else\textsubscript{\ensuremath{x}}\fi}
\newfontfamily\popdisplay{ArchivoBlack-Regular.ttf}[Path=../../../fonts/]
\newfontfamily\poplogo{SpaceGrotesk-Bold.ttf}[Path=../../../fonts/]
\newfontfamily\popsemi{PlusJakartaSans-SemiBold.ttf}[Path=../../../fonts/]
\newfontfamily\popxbold{PlusJakartaSans-ExtraBold.ttf}[Path=../../../fonts/]
\newfontfamily\popxboldls{PlusJakartaSans-ExtraBold.ttf}[Path=../../../fonts/,LetterSpace=3.0]

\setlist[enumerate]{leftmargin=1.7em,itemsep=5pt,topsep=4pt}
\setlist[itemize]{leftmargin=1.7em,itemsep=4pt,topsep=4pt,label={\color{poporange}\textbullet}}
\setlength{\parindent}{0pt}
\setlength{\parskip}{5pt}
\renewcommand{\arraystretch}{1.25}

% pgfplots : style Pop par défaut
\pgfplotsset{
  every axis/.append style={axis line style={popink,line width=1pt},tick style={popink},
    grid style={popline,line width=0.5pt},label style={font=\small},tick label style={font=\footnotesize}},
  popcurve/.style={poporange,line width=1.8pt,no markers,smooth},
}
% Même style disponible dans un \draw TikZ simple (hors pgfplots)
\tikzset{popcurve/.style={poporange,line width=1.8pt}}
\tikzset{popgrid/.style={popline,very thin}}
\colorlet{popcurve}{poporange} % le nom du style est parfois utilisé comme couleur

% ============================================================
% Pill / squiggle
% ============================================================
\newcommand{\poppill}[3]{%
  \tikz[baseline=-0.55ex]\node[draw=popink,line width=1.2pt,rounded corners=2.6pt,%
    fill=#2,inner xsep=6.5pt,inner ysep=3pt]{%
    \popxboldls\fontsize{7}{7}\selectfont\color{#3}\MakeUppercase{#1}};%
}
\newcommand{\popsquiggle}[1]{%
  \tikz[baseline=-0.4ex]{
    \draw[#1,line width=2.6pt,line cap=round]
      plot [smooth] coordinates {(0,0.09) (0.34,0.01) (0.68,0.21) (1.02,0.01) (1.36,0.10)};
  }%
}
% Mot-clé surligné (marqueur orange sous le texte)
\newcommand{\cle}[1]{{\popxbold\color{popdark}#1}}

% ============================================================
% En-tête / pied
% ============================================================
\pagestyle{fancy}
\fancyhf{}
\renewcommand{\headrulewidth}{0pt}
\renewcommand{\footrulewidth}{0pt}
\lhead{\raisebox{-0.15ex}{\poplogo\fontsize{13}{13}\selectfont\color{popink}OnBuch\color{poporange}+}}
\rhead{\poppill{\DOCMATIERE\ \textbullet\ \DOCNIVEAU\ \textbullet\ Leçon \DOCLECON}{white}{popink}}
\lfoot{\popxbold\fontsize{7.6}{7.6}\selectfont\color{popmuted}\MakeUppercase{Cours signé OnBuch+}\ \textbullet\ {\popsemi\DOCTITRE}}
\rfoot{\tikz[baseline=-0.6ex]\node[rounded corners=8.5pt,fill=popink,minimum height=17pt,minimum width=17pt,inner xsep=5pt,inner sep=0]{\popxbold\fontsize{8}{8}\selectfont\color{popcream}\thepage\,/\,\pageref*{LastPage}};}

% ============================================================
% Titres
% ============================================================
\newcommand{\chaptitre}[2]{%
  \begin{tcolorbox}[enhanced,colback=poporange,colframe=popink,boxrule=2.6pt,arc=11pt,
      drop shadow=popink,left=15pt,right=15pt,top=12pt,bottom=11pt,before skip=3pt,after skip=18pt]
    \poppill{#2}{popink}{popcream}\par\vspace{9pt}
    {\raggedright\hyphenpenalty=10000\exhyphenpenalty=10000\popdisplay\fontsize{22}{24}\selectfont\color{popcream}\MakeUppercase{#1}\par}
  \end{tcolorbox}
}
% Section numérotée : pastille orange avec le numéro + titre Archivo
\newcounter{popsec}
\newcommand{\coursec}[1]{%
  \stepcounter{popsec}\setcounter{popsub}{0}%
  \par\vspace{16pt}\noindent
  \begin{minipage}{\linewidth}
  \tikz[baseline=-0.7ex]\node[draw=popink,line width=1.6pt,fill=poporange,rounded corners=4pt,minimum size=22pt,inner sep=0]{\popdisplay\fontsize{12}{12}\selectfont\color{popcream}\thepopsec};%
  \hspace{8pt}{\popdisplay\fontsize{15}{17}\selectfont\color{popink}\MakeUppercase{#1}}\par
  \vspace{4pt}\noindent{\color{poporange}\rule{36pt}{3.6pt}}
  \end{minipage}\par\nopagebreak\vspace{8pt}
}
\newcounter{popsub}
\newcommand{\courssub}[1]{%
  \stepcounter{popsub}%
  \par\vspace{9pt}\noindent{\popxbold\fontsize{11.5}{13}\selectfont\color{popink}{\color{poporange}\thepopsec.\thepopsub}\hspace{6pt}#1}\par\nopagebreak\vspace{4pt}
}

% ============================================================
% Boîtes pédagogiques — même recette : fond blanc, bordure épaisse colorée,
% ombre dure encre, badge plein. Titre optionnel [#1].
% ============================================================
\tcbset{popbase/.style={breakable,enhanced,colback=white,boxrule=2.3pt,arc=9pt,
  shadow={3pt}{-3pt}{0pt}{popink},left=14pt,right=14pt,top=11pt,bottom=11pt,before skip=11pt,after skip=15pt,
  fontupper=\popsemi\fontsize{10.6}{14.5}\selectfont\color{popink},
  fontlower=\popsemi\fontsize{10.6}{14.5}\selectfont\color{popink}}}
% Coche dessinée (le glyphe ✓ n'existe pas dans les polices du document)
\newcommand{\popcheck}[1]{\tikz[baseline=-0.5ex]\draw[#1,line width=1.6pt,line cap=round,line join=round](-0.13,0.0)--(-0.04,-0.09)--(0.13,0.1);}
\providecommand{\coche}{\popcheck{popgreen}}
\providecommand{\resultat}[1]{\par\smallskip\noindent\fcolorbox{popgreen}{popgreenL}{\parbox{\dimexpr\linewidth-2\fboxsep-2\fboxrule\relax}{\textbf{Résultat.} #1}}\par\smallskip}
\newcommand{\popboxhead}[3]{\poppill{#1}{#2}{white}\ifx\relax#3\relax\else\hspace{6pt}{\popxbold\fontsize{10.6}{12}\selectfont\color{popink}#3}\fi\par\vspace{7pt}}

\newtcolorbox{aretenir}[1][]{popbase,colframe=popgreen,before upper={\popboxhead{À retenir}{popgreen}{#1}}}
% Texte en anglais (document de lecture, dialogue, extrait) : cadre
% d'étude. Un titre optionnel (source, auteur) s'affiche dans l'en-tête.
\newtcolorbox{textebox}[1][]{popbase,colframe=popblue,before upper={\popboxhead{Text}{popblue}{#1}\begin{otherlanguage}{english}},after upper={\end{otherlanguage}}}
% Marques ✓ / ✗ (forme correcte / fautive) souvent utilisées par les modèles
\providecommand{\cross}{\textcolor{poppink}{\ensuremath{\times}}}
\providecommand{\xmark}{\cross}\providecommand{\crossmark}{\cross}\providecommand{\cmark}{\ensuremath{\checkmark}}
% Ligne de réponse / justification à compléter (inventée par les modèles)
\providecommand{\justification}[1]{\par\noindent\textit{Justification~:} \dotfill\par}
% \textipa (package tipa non chargé) : la police ne couvre pas l'API -> simple passage
\providecommand{\textipa}[1]{#1}
% Soulignement (ulem non chargé) : \uline, \ul
\providecommand{\uline}[1]{\underline{#1}}\providecommand{\ul}[1]{\underline{#1}}
% \glossaire{\item ...} inventée par les modèles : liste de glossaire
\providecommand{\glossaire}[1]{\begin{itemize}#1\end{itemize}}
% \alert (beamer) inventée par les modèles : texte mis en évidence
\providecommand{\alert}[1]{\textbf{\textcolor{poppink}{#1}}}
% variantes ASCII d'ordinaux écrites en macro
\providecommand{\iere}{\textsuperscript{re}}\providecommand{\ieme}{\textsuperscript{e}}\providecommand{\ere}{\textsuperscript{re}}\providecommand{\eme}{\textsuperscript{e}}\providecommand{\er}{\textsuperscript{er}}\providecommand{\ie}{\textsuperscript{e}}
% Ordinaux français écrits en macro par les modèles : 1\ière, 3\ième
\newcommand{\ière}{\textsuperscript{re}}\newcommand{\ième}{\textsuperscript{e}}
% \exemple (inventée par les modèles) : étiquette « Exemple : »
\providecommand{\exemple}{\textbf{Exemple~:}\ }
% \square utilisable aussi hors mode math (cases à cocher)
\AtBeginDocument{\let\squareMath\square\renewcommand{\square}{\ensuremath{\squareMath}}}
% Mot ou phrase en anglais à l'intérieur d'un paragraphe français
\newcommand{\eng}[1]{\ifmmode\text{\textit{#1}}\else\begingroup\selectlanguage{english}\itshape #1\endgroup\fi}
\let\esp\eng % alias toléré
% flèches d'intonation inventées par les modèles
\providecommand{\textsearrow}{}\renewcommand{\textsearrow}{\ensuremath{\searrow}}\providecommand{\textnearrow}{}\renewcommand{\textnearrow}{\ensuremath{\nearrow}}
% \fr{..} : mot français (inventé par les modèles)
\providecommand{\fr}[1]{\textit{#1}}
\newtcolorbox{exemplebox}[1][]{popbase,colframe=poporange,before upper={\popboxhead{Exemple}{poporange}{#1}}}
\newtcolorbox{definition}[1][]{popbase,colframe=popblue,before upper={\popboxhead{Définition}{popblue}{#1}}}
\newtcolorbox{propriete}[1][]{popbase,colframe=poppurple,before upper={\popboxhead{Propriété}{poppurple}{#1}}}
\newtcolorbox{methode}[1][]{popbase,colframe=popgold,before upper={\popboxhead{Méthode}{popgold}{#1}}}
\newtcolorbox{attention}[1][]{popbase,colframe=poppink,before upper={\popboxhead{Attention}{poppink}{#1}}}
\newtcolorbox{experience}[1][]{popbase,colframe=popblue,colback=popblueL,before upper={\popboxhead{Expérience}{popblue}{#1}}}
% « Le savais-tu ? » : carte violette pleine (contexte, culture, Cameroun)
\newtcolorbox{savaistu}[1][]{popbase,colback=poppurple,colframe=popink,
  fontupper=\popsemi\fontsize{10.4}{14.2}\selectfont\color{white},
  before upper={\poppill{Le savais-tu\,?}{white}{poppurple}\ifx\relax#1\relax\else\hspace{6pt}{\popxbold\color{white}#1}\fi\par\vspace{7pt}}}
% Exercice résolu : énoncé en haut, \tcblower puis la solution sur fond crème
\newcounter{popexo}\newcounter{popexr}
\newtcolorbox{exoresolu}[1][]{popbase,colframe=poporange,colbacklower=popcream,
  segmentation style={popink,line width=1.2pt,dashed},
  before upper={\stepcounter{popexr}\popboxhead{Exercice résolu \thepopexr}{poporange}{#1}},
  before lower={\poppill{Solution}{popink}{popcream}\par\vspace{6pt}}}
% Exercice (non résolu en place) — la correction est dans la partie Corrigés
\newtcolorbox{exercice}[1][]{popbase,colframe=popink,
  before upper={\stepcounter{popexo}\popboxhead{Exercice \thepopexo}{popink}{#1}}}
% Corrigé
\newtcolorbox{corrige}[1][]{popbase,colframe=popgreen,colback=popgreenL,
  before upper={\popboxhead{Corrigé}{popgreen}{#1}}}
% Objectifs / prérequis / bilan
\newtcolorbox{objectifs}{popbase,colframe=popink,colback=poporangeL,
  before upper={\popboxhead{Objectifs de la leçon}{popink}{}}}
\newtcolorbox{prerequis}{popbase,colframe=popink,colback=popblueL,
  before upper={\popboxhead{Prérequis}{popblue}{}}}
\newtcolorbox{bilan}{popbase,colframe=popink,colback=white,boxrule=2.8pt,
  before upper={\popboxhead{Fiche bilan}{poporange}{L'essentiel à connaître pour l'examen}}}
% Figure encadrée avec légende
\newtcolorbox{popfigure}[1][]{enhanced,breakable=false,colback=white,colframe=popline,boxrule=1.4pt,arc=8pt,
  left=10pt,right=10pt,top=10pt,bottom=8pt,before skip=10pt,after skip=12pt,halign=center,
  lower separated=false,halign lower=center,
  fontlower=\popsemi\fontsize{9}{11.5}\selectfont\color{popmuted},
  #1}
\newcommand{\legende}[1]{\par\vspace{6pt}{\popsemi\fontsize{9}{11.5}\selectfont\color{popmuted}#1}}

% ============================================================
% Signature OnBuch+
% ============================================================
\newcommand{\poplogoinline}[1]{{\poplogo\fontsize{#1}{#1}\selectfont\color{popink}OnBuch\color{poporange}+}}
\newcommand{\signatureonbuch}{%
  \begin{tcolorbox}[enhanced,colback=popink,colframe=popink,boxrule=2.2pt,arc=10pt,
      drop shadow=poporange,left=14pt,right=14pt,top=10pt,bottom=10pt,before skip=0pt,after skip=0pt]
    \tikz[baseline=-0.7ex]\node[circle,fill=popgreen,draw=popcream,line width=1.3pt,minimum size=22pt,inner sep=0]{\popcheck{white}};%
    \hspace{9pt}\begin{minipage}[c]{0.62\linewidth}
      {\popxbold\fontsize{12}{13}\selectfont\color{popcream}Signé\ \textbullet\ {\poplogo OnBuch\color{poporange}+}}\par\vspace{2pt}
      {\raggedright\popsemi\fontsize{8}{10.5}\selectfont\color{popline}\MakeUppercase{Équipe pédagogique OnBuch+}\\\MakeUppercase{Conforme au programme officiel MINESEC}\par}
    \end{minipage}\hfill
    {\poplogo\fontsize{22}{22}\selectfont\color{popcream}OnBuch\color{poporange}+}
  \end{tcolorbox}%
}

% ============================================================
% Page de garde
% #1 titre · #2 matière · #3 niveau · #4 module · #5 n° leçon · #6 durée ·
% #7 description / objectifs (texte ou itemize)
% ============================================================
\newcommand{\pagegarde}[7]{%
  \thispagestyle{empty}
  \newgeometry{a4paper,margin=1.7cm,top=1.6cm,bottom=1.4cm}%
  \begin{tikzpicture}[remember picture,overlay]
    \fill[poporange!18] ([xshift=-3.2cm,yshift=-3.6cm]current page.north east) circle (4.6cm);
    \fill[poppurple!14] ([xshift=2.4cm,yshift=7.6cm]current page.south west) circle (3.8cm);
  \end{tikzpicture}%
  {\poplogo\fontsize{30}{30}\selectfont\color{popink}OnBuch\color{poporange}+}\hfill
  \raisebox{0.6ex}{\poppill{Cours signé \textbullet\ Premium}{popink}{popcream}}\par
  \vspace{1pt}\popsquiggle{poppurple}\par
  \vspace{8pt}
  \begin{tcolorbox}[enhanced,colback=poporange,colframe=popink,boxrule=2.8pt,arc=13pt,
      drop shadow=popink,left=18pt,right=18pt,top=12pt,bottom=13pt,after skip=10pt]
    \poppill{#2 \textbullet\ #3}{popink}{popcream}\hspace{5pt}\poppill{#4}{white}{popink}\par\vspace{8pt}
    {\popdisplay\fontsize{14}{14}\selectfont\color{popink}LEÇON #5}\par\vspace{4pt}
    {\raggedright\hyphenpenalty=10000\exhyphenpenalty=10000\popdisplay\fontsize{25}{27}\selectfont\color{popcream}\MakeUppercase{#1}\par}
    \vspace{8pt}
    {\popxbold\fontsize{9.5}{9.5}\selectfont\color{popink}\MakeUppercase{Durée conseillée : #6}}
  \end{tcolorbox}
  \begin{tcolorbox}[enhanced,colback=white,colframe=popink,boxrule=2.2pt,arc=10pt,
      drop shadow=popink,left=16pt,right=16pt,top=10pt,bottom=10pt,after skip=8pt]
    \poppill{Dans cette leçon}{poppurple}{white}\par\vspace{5pt}
    \popsemi\fontsize{9.3}{12.6}\selectfont\color{popink}#7
  \end{tcolorbox}
  \noindent\begin{minipage}[t]{0.487\linewidth}
    \begin{tcolorbox}[enhanced,colback=popgreen,colframe=popink,boxrule=2.2pt,arc=10pt,
        drop shadow=popink,left=11pt,right=11pt,top=10pt,bottom=10pt,height=3.1cm,valign=center]
      \tcbox[colback=white,colframe=popink,boxrule=1.3pt,arc=5pt,boxsep=0pt,nobeforeafter,
        left=3pt,right=3pt,top=3pt,bottom=3pt,valign=center]{\qrcode[height=1.9cm]{https://play.google.com/store/apps/details?id=com.luvvix.onbuch&hl=fr}}%
      \hspace{8pt}\begin{minipage}[c]{0.5\linewidth}
        {\popdisplay\fontsize{11}{12.5}\selectfont\color{white}\MakeUppercase{Télécharge OnBuch}}\par\vspace{4pt}
        {\raggedright\popsemi\fontsize{8.4}{11}\selectfont\color{white}Gratuit \textbullet\ Google Play\\Tuteur IA, annales, concours, résultats\par}
      \end{minipage}
    \end{tcolorbox}
  \end{minipage}\hfill
  \begin{minipage}[t]{0.487\linewidth}
    \begin{tcolorbox}[enhanced,colback=poppurple,colframe=popink,boxrule=2.2pt,arc=10pt,
        drop shadow=popink,left=11pt,right=11pt,top=10pt,bottom=10pt,height=3.1cm,valign=center]
      \tikz[baseline=-0.7ex]\node[circle,fill=white,draw=popink,line width=1.3pt,minimum size=1.6cm,inner sep=0]{\popdisplay\fontsize{24}{24}\selectfont\color{poppurple}+};%
      \hspace{8pt}\begin{minipage}[c]{0.6\linewidth}
        {\popdisplay\fontsize{11}{12.5}\selectfont\color{white}\MakeUppercase{Passe à OnBuch+}}\par\vspace{4pt}
        {\raggedright\popsemi\fontsize{8.4}{11}\selectfont\color{white}Tous les cours signés, Tuteur IA illimité, fiches PDF — onglet OnBuch+ de l'app\par}
      \end{minipage}
    \end{tcolorbox}
  \end{minipage}
  \vspace{8pt}
  \popcontenu
  \vfill
  \signatureonbuch
  \restoregeometry
  \newpage
}

% Rangée « Ce cours contient » (page de garde)
\newcommand{\poptile}[3]{% #1 couleur #2 gros texte #3 légende
  \begin{tcolorbox}[enhanced,colback=#1,colframe=popink,boxrule=2pt,arc=9pt,shadow={2.5pt}{-2.5pt}{0pt}{popink},
      left=6pt,right=6pt,top=9pt,bottom=9pt,halign=center,valign=center,height=1.75cm,width=\linewidth,nobeforeafter]
    {\popdisplay\fontsize{12.5}{13}\selectfont\color{white}\MakeUppercase{#2}}\par\vspace{3pt}
    {\centering\popsemi\fontsize{7.8}{9.5}\selectfont\color{white}#3\par}
  \end{tcolorbox}}
\newcommand{\popcontenu}{%
  \par\noindent{\popxboldls\fontsize{7.5}{7.5}\selectfont\color{popmuted}CE COURS SIGNÉ CONTIENT}\par\vspace{7pt}
  \noindent\begin{minipage}[t]{0.235\linewidth}\poptile{popblue}{Cours}{complet, illustré, pas à pas}\end{minipage}\hfill
  \begin{minipage}[t]{0.235\linewidth}\poptile{poporange}{Méthodes}{et exercices résolus}\end{minipage}\hfill
  \begin{minipage}[t]{0.235\linewidth}\poptile{poppink}{Exercices}{type examen + corrigés}\end{minipage}\hfill
  \begin{minipage}[t]{0.235\linewidth}\poptile{popgreen}{Fiche bilan}{l'essentiel à retenir}\end{minipage}\par}

% Sommaire
\newtcolorbox{sommaire}{enhanced,colback=white,colframe=popink,boxrule=2.2pt,arc=10pt,
  drop shadow=popink,left=18pt,right=18pt,top=13pt,bottom=13pt,before skip=6pt,after skip=20pt,
  before upper={\poppill{Sommaire}{poppurple}{white}\par\vspace{9pt}\popsemi\fontsize{10.6}{14.5}\selectfont\color{popink}}}

% Signature de fin de cours — poussée tout en bas de la dernière page
\newcommand{\findecours}{%
  \par\vfill\nopagebreak
  \begin{center}\popsquiggle{poporange}\end{center}\vspace{4pt}
  \signatureonbuch
}
\AtBeginDocument{\def\llbracket{\mathopen{[\mkern-3.5mu[}}\def\rrbracket{\mathclose{]\mkern-3.5mu]}}} % intervalles d'entiers (glyphe ⟦ absent de la police)
% Glyphes absents de Fira Math : redéfinis à partir de glyphes disponibles
\AtBeginDocument{%
  \def\setminus{\mathbin{\backslash}}\let\smallsetminus\setminus
  \def\vdots{\mathord{\vbox{\baselineskip3.2pt\lineskiplimit0pt\kern4pt\hbox{.}\hbox{.}\hbox{.}}}}%
  \def\ddots{\mathinner{\mkern1mu\raise6.5pt\hbox{.}\mkern2mu\raise3.6pt\hbox{.}\mkern2mu\raise0.7pt\hbox{.}\mkern1mu}}%
  \def\triangle{\mathord{\scalebox{0.9}{$\symup{\Delta}$}}}%
  \def\nabla{\mathord{\raisebox{\depth}{\scalebox{1}[-1]{$\symup{\Delta}$}}}}%
  \def\checkmark{\popcheck{popdark}}%
  \def\star{\mathbin{\ast}}\def\bigstar{\mathord{\ast}}%
  \def\diamond{\mathbin{\scalebox{0.6}{$\Diamond$}}}%
  \def\bigcup{\mathop{\mathchoice{\raisebox{-0.3em}{\scalebox{1.7}{$\cup$}}}{\scalebox{1.25}{$\cup$}}{\cup}{\cup}}\displaylimits}%
  \def\bigcap{\mathop{\mathchoice{\raisebox{-0.3em}{\scalebox{1.7}{$\cap$}}}{\scalebox{1.25}{$\cap$}}{\cap}{\cap}}\displaylimits}%
  \def\lesseqqgtr{\mathrel{\lessgtr}}\def\bigoplus{\mathop{\oplus}\displaylimits}%
}
\providecommand{\ssi}{si et seulement si} % abréviation fréquente
\AtBeginDocument{\providecommand{\wideparen}{\overparen}} % arc de cercle

\begin{document}

\pagegarde{Lesson 59 — Vocabulary: Words and expressions related to difficulties with the use of ICTs at work or for learning}{Anglais}{Tle SH}{Chapter 5 : Media and Communication}{EN59}{2 h}{Cette leçon de vocabulaire présente les mots et expressions anglais permettant de décrire les difficultés liées à l'usage des TIC au travail et dans les apprentissages (pannes, bugs, problèmes de connexion, difficultés d'utilisation). Elle insiste sur les collocations, les faux amis et le réemploi du lexique dans des contextes authentiques camerounais et anglophones.
\vspace{4pt}
\begin{multicols}{2}\small
\begin{itemize}
\item À la fin de cette leçon, je sais nommer en anglais les principales difficultés techniques liées aux TIC (pannes, bugs, lenteurs, coupures).
\item À la fin de cette leçon, je sais utiliser les collocations correctes avec les noms et verbes du domaine (to crash, to log in, to troubleshoot...).
\item À la fin de cette leçon, je sais employer des expressions idiomatiques pour décrire un problème informatique (the system is down, to be at a loss...).
\item À la fin de cette leçon, je sais éviter les faux amis et les calques du français (deception, actually, to assist...).
\item À la fin de cette leçon, je sais former des mots de la même famille par préfixes et suffixes (connect, connection, connectivity, disconnect...).
\item À la fin de cette leçon, je sais prononcer et orthographier correctement le lexique étudié.
\end{itemize}\end{multicols}}

\begin{sommaire}
\begin{multicols}{2}
\begin{enumerate}
\item Mise en situation et découverte du lexique
\item Champs lexicaux : les difficultés avec les TIC
\item Collocations et expressions idiomatiques
\item Faux amis et calques du français
\item Familles de mots, prononciation et orthographe
\item Réemploi en contexte et production guidée
\item Activité d'intégration
\item Méthodes et exercices résolus
\item Exercices
\item Corrigés des exercices
\item Fiche bilan
\end{enumerate}
\end{multicols}
\end{sommaire}

\begin{objectifs}
\begin{itemize}
\item À la fin de cette leçon, je sais nommer en anglais les principales difficultés techniques liées aux TIC (pannes, bugs, lenteurs, coupures).
\item À la fin de cette leçon, je sais utiliser les collocations correctes avec les noms et verbes du domaine (to crash, to log in, to troubleshoot...).
\item À la fin de cette leçon, je sais employer des expressions idiomatiques pour décrire un problème informatique (the system is down, to be at a loss...).
\item À la fin de cette leçon, je sais éviter les faux amis et les calques du français (deception, actually, to assist...).
\item À la fin de cette leçon, je sais former des mots de la même famille par préfixes et suffixes (connect, connection, connectivity, disconnect...).
\item À la fin de cette leçon, je sais prononcer et orthographier correctement le lexique étudié.
\item À la fin de cette leçon, je sais réemployer ce vocabulaire dans des phrases, dialogues et courts textes à l'oral comme à l'écrit.
\item À la fin de cette leçon, je sais raconter une difficulté rencontrée avec les TIC et proposer une solution.
\end{itemize}
\end{objectifs}

\begin{prerequis}
\begin{itemize}
\item Vocabulaire de base des TIC vu en Première (computer, internet, email, website, password).
\item Le présent simple et le present continuous pour décrire des situations habituelles et en cours.
\item Le prétérit simple pour raconter un incident passé.
\item Les modaux can / can't pour exprimer la capacité ou l'impossibilité.
\item La formation des mots par préfixes et suffixes (notion vue en Première).
\end{itemize}
\end{prerequis}

\newpage

\coursec{Mise en situation et découverte du lexique}

\begin{textebox}[A real-life situation]
At the Ministry of Secondary Education in Yaoundé, a secretary stares at her frozen screen just minutes before sending an urgent report. In Bamenda, a student loses his online exam answers because the connection suddenly drops. “The system is down again!”, “My laptop crashed!”, “I can’t log in!” — these complaints are heard every day in Cameroonian offices, schools and cybercafés. Knowing how to name and describe these ICT problems in correct English is now an essential skill for work, study and the Bac exam.
\end{textebox}

\textbf{Problématique.} Comment nommer et décrire en anglais correct les difficultés que l’on rencontre avec les TIC (ordinateurs, internet, logiciels) au travail ou à l’école ? Pour y répondre, nous allons d’abord lire un récit authentique, repérer les mots nouveaux, deviner leur sens grâce au contexte, vérifier nos hypothèses, puis classer ce vocabulaire en catégories.

\courssub{Lecture du texte support : \eng{A Bad Day at the Cybercafé}}

Lisons le récit de Ndi, un étudiant de l’Université de Buea, qui a vécu une très mauvaise journée à cause des technologies. Lis le texte une première fois \textbf{sans dictionnaire} : essaie seulement de comprendre l’histoire globale.

\begin{textebox}[A Bad Day at the Cybercafé — by Ndi, a student in Buea]
Yesterday, I was typing my end-of-year project at a cybercafé in Buea when the screen suddenly froze. I pressed every key, but nothing moved. I tried to restart the computer, but it crashed completely and made a strange noise. The owner of the cybercafé told me there was probably a bug in the program I was using. To make matters worse, the internet connection was so slow that I could not even log in to my email account to retrieve my backup file. I waited for twenty minutes, but the page would not load. Then, when the power cut came, all the computers went off and I lost everything — three hours of hard work! I felt helpless and frustrated. The owner apologised and said the breakdown was caused by the heavy rain. Next time, he promised, he would install an antivirus and save the customers’ files on an external hard drive. As for me, I have learnt my lesson: I will always save my work every five minutes and keep a copy on a flash drive.
\end{textebox}

\begin{definition}[Glossaire du texte]
\begin{itemize}
\item \cle{to type} (verbe) : taper (au clavier) — « taïp »
\item \cle{screen} (nom) : écran — « skrine »
\item \cle{to freeze / froze / frozen} (verbe irrégulier) : se figer, se bloquer — « frize / froze / fro-zen »
\item \cle{to restart} (verbe) : redémarrer
\item \cle{to crash} (verbe) : planter (ordinateur, programme) — « krach »
\item \cle{bug} (nom) : bogue, erreur dans un programme — « beug »
\item \cle{slow connection} (nom) : connexion lente
\item \cle{to log in} (verbe) : se connecter (à un compte)
\item \cle{backup file} (nom) : fichier de sauvegarde
\item \cle{power cut} (nom) : coupure de courant, délestage
\item \cle{breakdown} (nom) : panne — « breik-daoun »
\item \cle{helpless} (adjectif) : impuissant, désemparé
\item \cle{flash drive} (nom) : clé USB
\end{itemize}
\end{definition}

\textbf{Questions de compréhension globale (oral).}
\begin{enumerate}
\item Where was Ndi when the problem started?
\item What was he doing when the screen froze?
\item How many different problems did he face? Name them.
\item What lesson did he learn at the end?
\end{enumerate}

\begin{corrige}[Questions de compréhension]
\begin{enumerate}
\item He was at a cybercafé in Buea. « Il était dans un cybercafé à Buea. »
\item He was typing his end-of-year project. « Il tapait son projet de fin d’année. »
\item He faced four main problems: the screen froze, the computer crashed, the connection was too slow, and there was a power cut. « Il a connu quatre problèmes principaux : l’écran s’est figé, l’ordinateur a planté, la connexion était trop lente et il y a eu une coupure de courant. »
\item He learnt that he must always save his work regularly and keep a copy on a flash drive. « Il a appris qu’il doit toujours sauvegarder son travail régulièrement et en garder une copie sur une clé USB. »
\end{enumerate}
\end{corrige}

\courssub{Repérage des mots nouveaux liés aux difficultés}

Relisons le texte et soulignons tous les mots et expressions qui désignent un \textbf{problème technique}. Nous obtenons six mots-clés :

\begin{center}
\begin{tabularx}{\linewidth}{|l|X|X|}\hline
\rowcolor{popblueL} \textbf{Mot anglais} & \textbf{Nature} & \textbf{Traduction française} \\ \hline
\cle{freeze} (froze, frozen) & verbe irrégulier & se figer, se bloquer (écran) \\ \hline
\cle{crash} & verbe et nom & planter ; un plantage \\ \hline
\cle{bug} & nom & bogue, erreur de programme \\ \hline
\cle{breakdown} & nom & panne (générale) \\ \hline
\cle{power cut} & nom & coupure de courant, délestage \\ \hline
\cle{slow connection} & nom & connexion (internet) lente \\ \hline
\end{tabularx}
\end{center}

\begin{exemplebox}[Exemples tirés du texte, traduits]
\eng{The screen froze.} « L’écran s’est figé / s’est bloqué. »

\eng{The computer crashed.} « L’ordinateur a planté. »

\eng{There was a bug in the program.} « Il y avait une bogue dans le programme. »

\eng{When the power cut came, I lost everything.} « Quand la coupure de courant est arrivée, j’ai tout perdu. »

\eng{The connection was so slow that I could not log in.} « La connexion était si lente que je ne pouvais même pas me connecter. »
\end{exemplebox}

\begin{propriete}[Le verbe irrégulier \eng{to freeze}]
\eng{To freeze} est un verbe irrégulier : \textbf{freeze — froze — frozen}. Attention à ne pas conjuguer \emph{freezed}, qui n’existe pas.

\begin{center}
\begin{tabular}{|l|l|l|}\hline
\rowcolor{popblueL} \textbf{Base verbale} & \textbf{Prétérit} & \textbf{Participe passé} \\ \hline
freeze & froze & frozen \\ \hline
\end{tabular}
\end{center}

\eng{My computer has frozen again.} « Mon ordinateur s’est encore bloqué. » (present perfect : has + frozen)
\end{propriete}

\courssub{Stratégie : deviner le sens d’un mot à partir du contexte}

Avant d’ouvrir le dictionnaire, un bon lecteur devine le sens des mots nouveaux. C’est une compétence essentielle pour l’épreuve de \emph{reading} au Bac : tu n’auras pas de dictionnaire !

\begin{methode}[Comment deviner le sens d’un mot inconnu]
\begin{enumerate}
\item \textbf{Identifier sa nature} : est-ce un nom, un verbe, un adjectif ? Observe sa place dans la phrase et ses terminaisons (\emph{-ed} pour un verbe au passé, \emph{a/an/the} devant un nom).
\item \textbf{Observer les mots qui l’entourent} : le contexte donne des indices (cause, conséquence, synonyme, opposition).
\item \textbf{Formuler une hypothèse} : propose une traduction possible en français, même approximative.
\item \textbf{Vérifier} : contrôle ton hypothèse avec le glossaire ou le dictionnaire, puis relis la phrase pour confirmer.
\end{enumerate}
\end{methode}

\begin{exoresolu}[Devinette lexicale guidée]
Énoncé. Sans dictionnaire, devine le sens de \eng{froze} dans la phrase : \eng{I was typing my project when the screen suddenly froze. I pressed every key, but nothing moved.}
\tcblower
Solution détaillée.
\begin{enumerate}
\item \textbf{Nature} : \eng{froze} est placé après le sujet \eng{the screen} ; c’est un verbe au prétérit (forme irrégulière, sans complément).
\item \textbf{Contexte} : la phrase suivante dit \eng{I pressed every key, but nothing moved} (« j’ai appuyé sur toutes les touches, mais rien ne bougeait »). L’écran ne réagit plus.
\item \textbf{Hypothèse} : \eng{froze} signifie probablement « s’est bloqué / s’est figé ».
\item \textbf{Vérification} : le glossaire confirme : \eng{to freeze, froze, frozen} = se figer, se bloquer. Hypothèse correcte !
\end{enumerate}
\end{exoresolu}

\begin{attention}[Piège : l’ordre des mots avec \eng{everything}]
Ne traduis jamais \eng{I lost everything} par « j’ai perdu tout » mot à mot, et surtout ne dis pas \emph{I lost all} en calquant « j’ai tout perdu ». En anglais, les pronoms indéfinis comme \eng{everything}, \eng{nothing}, \eng{something} se placent \textbf{après le verbe} :

\eng{I lost everything.} « J’ai tout perdu. »

\eng{She said nothing.} « Elle n’a rien dit. »

De même, attention à \eng{to lose} (perdre, prononcé « louze ») qui ne s’écrit jamais \emph{to loose} : \eng{loose} (« louce ») est un adjectif qui signifie « lâche, desserré ».
\end{attention}

\courssub{Vérification des hypothèses et première classification}

Après avoir vérifié nos hypothèses avec le glossaire, classons oralement les six mots-clés en quatre catégories : les problèmes \textbf{matériels} (hardware), \textbf{logiciels} (software), de \textbf{réseau} (network) et \textbf{humains} (user problems).

\begin{center}
\begin{tabularx}{\linewidth}{|l|X|X|}\hline
\rowcolor{popblueL} \textbf{Catégorie} & \textbf{Mots du texte} & \textbf{Explication} \\ \hline
Hardware (matériel) & power cut, breakdown & l’électricité ou la machine elle-même tombe en panne \\ \hline
Software (logiciel) & bug, crash, freeze & le programme ou l’ordinateur cesse de fonctionner \\ \hline
Network (réseau) & slow connection & internet est lent ou coupé \\ \hline
User problems (humain) & (I lost everything) & l’utilisateur n’a pas sauvegardé son travail \\ \hline
\end{tabularx}
\end{center}

\begin{popfigure}
\begin{tikzpicture}[mindmap, grow cyclic, every node/.style={concept, align=center, font=\small},
root concept/.append style={concept color=popblue, font=\bfseries},
level 1/.append style={level distance=3.4cm, sibling angle=90}]
\node[root concept]{ICT difficulties}
  child[concept color=poporange]{ node[align=center]{Hardware\\ power cut\\ breakdown} }
  child[concept color=popgreen]{ node[align=center]{Software\\ bug\\ crash\\ freeze} }
  child[concept color=poppurple]{ node[align=center]{Network\\ slow connection\\ no signal} }
  child[concept color=poppink]{ node[align=center]{User problems\\ forgot to save\\ wrong password} };
\end{tikzpicture}
\legende{Carte mentale : les quatre familles de difficultés liées aux TIC.}
\end{popfigure}

\begin{exercice}[À toi de jouer — classification orale]
En binôme, classe les expressions suivantes dans les quatre catégories (hardware, software, network, user problems), puis utilise chacune dans une phrase complète : \eng{a dead battery}, \eng{a virus}, \eng{no Wi-Fi signal}, \eng{a forgotten password}, \eng{a broken keyboard}, \eng{the system is down}.
\end{exercice}

\begin{corrige}[Exercice — classification]
\begin{itemize}
\item \textbf{Hardware} : a dead battery (batterie vide), a broken keyboard (clavier cassé).
\item \textbf{Software} : a virus, the system is down (le système est en panne).
\item \textbf{Network} : no Wi-Fi signal (pas de signal Wi-Fi).
\item \textbf{User problems} : a forgotten password (mot de passe oublié).
\end{itemize}
Exemple de phrase : \eng{I could not submit my homework because there was no Wi-Fi signal in my quarter.} « Je n’ai pas pu rendre mes devoirs parce qu’il n’y avait pas de signal Wi-Fi dans mon quartier. »
\end{corrige}

\begin{savaistu}[D’où vient le mot \eng{bug} ?]
Le mot \eng{bug} signifie littéralement « insecte » en anglais. Selon une histoire célèbre de l’informatique, en 1947, des techniciens américains ont découvert qu’un vrai papillon de nuit (a moth) coincé dans un ordinateur en provoquait la panne : depuis, on appelle \eng{bug} toute erreur dans un programme, et \eng{to debug} signifie « corriger les erreurs d’un programme ». Au Cameroun, le mot « bug » est d’ailleurs passé dans le langage courant : on dit souvent « le réseau bugue » pour parler d’une connexion instable !
\end{savaistu}

\begin{aretenir}[Ce qu’il faut retenir de cette étape]
\begin{itemize}
\item Six mots-clés pour décrire les pannes : \cle{freeze}, \cle{crash}, \cle{bug}, \cle{breakdown}, \cle{power cut}, \cle{slow connection}.
\item Pour deviner un mot inconnu : nature $\rightarrow$ contexte $\rightarrow$ hypothèse $\rightarrow$ vérification.
\item Les difficultés TIC se classent en quatre familles : \eng{hardware}, \eng{software}, \eng{network}, \eng{user problems}.
\item Ordre des mots : \eng{I lost everything} (jamais \emph{I lost all}).
\item Verbe irrégulier : \eng{freeze — froze — frozen} (jamais \emph{freezed}).
\end{itemize}
\end{aretenir}

\coursec{Champs lexicaux : les difficultés avec les TIC}

Dans la section précédente, nous avons découvert, à travers un dialogue et un court article, que les technologies de l'information et de la communication (ICTs) ne fonctionnent pas toujours comme prévu : l'ordinateur de Mme Ngo Bell a refusé de démarrer le matin de sa présentation, et les élèves du lycée de Bafoussam n'ont pas pu suivre leur cours en ligne. Pour parler de ces situations avec précision, il ne suffit pas de connaître quelques mots isolés : il faut maîtriser des \cle{champs lexicaux} organisés. Un champ lexical est un ensemble de mots et d'expressions qui se rapportent à un même thème. Dans cette section, nous allons explorer six champs lexicaux complémentaires : les pannes matérielles, les problèmes logiciels, les problèmes de réseau, les difficultés des utilisateurs, les conséquences et émotions, et enfin les solutions. À la fin, vous serez capable de décrire n'importe quel problème informatique et de proposer une solution, à l'oral comme à l'écrit.

\courssub{Observation : un texte pour découvrir les champs lexicaux}

Lisons d'abord ce court article original. Soulignez mentalement tous les mots et expressions qui désignent un problème ou une solution.

\begin{textebox}[A difficult Monday at the office]
Last Monday was a nightmare for the staff of a microfinance company in Douala. When the secretary switched on her computer at 7.30 a.m., nothing happened: the battery was dead and the charger was faulty. Her colleague's laptop was not better off — the screen was broken and the keyboard was not responding. At 9 a.m., the manager discovered that a virus had corrupted several files, and the whole system crashed twice before lunch. To make matters worse, the Internet connection was so slow that nobody could send e-mails, and by noon the network was completely down. “We are offline again!” complained the accountant, who had forgotten her password and got locked out of the accounting software. She felt helpless and frustrated: three hours of work had been lost because nobody had backed up the data. Finally, the company called a technician. He restarted the servers, installed an antivirus program, updated the software and reset the accountant's password. “Always back up your files,” he said before leaving, “and call our helpline as soon as something goes wrong.” By 4 p.m., everything was working again, but the day had been stressful and terribly time-consuming.
\end{textebox}

\textbf{Glossaire :}

\begin{center}
\begin{tabularx}{\linewidth}{|l|X|}\hline
\rowcolor{popblueL} \textbf{Mot anglais} & \textbf{Traduction française} \\ \hline
faulty (adj.) & défectueux \\ \hline
to crash (v.) & planter (pour un ordinateur) \\ \hline
down (adj.) & en panne (pour un réseau) \\ \hline
to get locked out (v.) & se retrouver bloqué, ne plus pouvoir accéder \\ \hline
to back up (v.) & sauvegarder \\ \hline
helpline (n.) & assistance téléphonique \\ \hline
time-consuming (adj.) & chronophage, qui prend du temps \\ \hline
\end{tabularx}
\end{center}

\textbf{Questions de compréhension :}
\begin{enumerate}
\item What was wrong with the secretary's computer?
\item Why could nobody send e-mails at 9 a.m.?
\item What did the technician do to solve the problems?
\end{enumerate}

Ce texte contient presque tout le lexique de la leçon. Classons-le maintenant méthodiquement.

\courssub{Champ 1 — Les pannes matérielles (hardware problems)}

Le \cle{hardware} désigne la partie matérielle de l'ordinateur : l'écran, le clavier, la batterie, l'imprimante. Attention : \eng{hardware} est \textbf{indénombrable} en anglais — on ne dit jamais \eng{hardwares}.

\begin{center}
\begin{tabularx}{\linewidth}{|X|X|X|}\hline
\rowcolor{popblueL} \textbf{English} & \textbf{Français} & \textbf{Exemple} \\ \hline
a breakdown & une panne & \eng{We had a breakdown just before the exam.} \\ \hline
a failure & une défaillance & \eng{a power failure} (une coupure de courant) \\ \hline
a fault & un défaut & \eng{There is a fault in the charger.} \\ \hline
out of order & en panne, hors service & \eng{The printer is out of order.} \\ \hline
a dead battery & une batterie à plat & \eng{My phone has a dead battery.} \\ \hline
a broken screen & un écran cassé & \eng{He dropped his laptop and now the screen is broken.} \\ \hline
the keyboard is not responding & le clavier ne répond pas & \eng{My keyboard is not responding.} \\ \hline
\end{tabularx}
\end{center}

\begin{attention}[Out of order : pour les machines seulement]
\eng{Out of order} s'emploie pour une machine en panne : \eng{The printer is out of order.} (« L'imprimante est en panne. »). On ne dit \textbf{pas} \eng{The secretary is out of order} pour dire qu'une personne est fatiguée ou malade ! Pour une personne, utilisez \eng{she is not feeling well} ou \eng{she is exhausted}. Notez aussi que \eng{out of order} ne se traduit pas ici par « hors d'usage » au sens juridique, mais simplement par « en panne ».
\end{attention}

\courssub{Champ 2 — Les problèmes logiciels (software problems)}

Le \cle{software} désigne les programmes (lui aussi indénombrable : jamais \eng{softwares}).

\begin{definition}[Bug]
A \cle{bug} is a small fault in a computer program that makes it work incorrectly. En français : un bogue, un défaut dans un programme. Exemple : \eng{There is a bug in the new version of the app.} (« Il y a un bogue dans la nouvelle version de l'application. »).
\end{definition}

\begin{center}
\begin{tabularx}{\linewidth}{|X|X|}\hline
\rowcolor{popblueL} \textbf{English} & \textbf{Français} \\ \hline
a bug & un bogue \\ \hline
a virus & un virus \\ \hline
to crash & planter \\ \hline
to freeze & se figer, se bloquer (l'écran) \\ \hline
an error message & un message d'erreur \\ \hline
a corrupted file & un fichier corrompu \\ \hline
to install / to uninstall & installer / désinstaller \\ \hline
to update & mettre à jour \\ \hline
\end{tabularx}
\end{center}

\eng{The system crashed and an error message appeared.} (« Le système a planté et un message d'erreur est apparu. ») — \eng{You should update your antivirus regularly.} (« Tu devrais mettre à jour ton antivirus régulièrement. »)

\begin{savaistu}[D'où vient le mot « bug » ?]
Le mot \eng{bug} (« insecte » en anglais courant) viendrait d'un vrai insecte ! En 1947, les ingénieurs de l'université Harvard auraient trouvé un papillon de nuit coincé dans un ordinateur, qui provoquait des pannes. Ils l'auraient collé dans leur registre avec la mention “first actual case of bug being found” (“premier cas réel d'insecte découvert”). Depuis, un défaut informatique s'appelle un \eng{bug}, et corriger un programme se dit \eng{to debug} (“déboguer”).
\end{savaistu}

\courssub{Champ 3 — Les problèmes de réseau (network problems)}

Au Cameroun comme ailleurs, beaucoup de difficultés viennent de la connexion elle-même.

\begin{center}
\begin{tabularx}{\linewidth}{|X|X|}\hline
\rowcolor{popblueL} \textbf{English} & \textbf{Français} \\ \hline
a slow connection & une connexion lente \\ \hline
no signal & pas de signal \\ \hline
a network failure & une panne de réseau \\ \hline
to disconnect & se déconnecter \\ \hline
to be offline & être hors ligne \\ \hline
bandwidth & la bande passante \\ \hline
a Wi-Fi hotspot & un point d'accès Wi-Fi \\ \hline
the network is down & le réseau est en panne \\ \hline
\end{tabularx}
\end{center}

\begin{exemplebox}[Exemples traduits]
\eng{The network is down.} = « Le réseau est en panne. »

\eng{There is no signal in the village, so I cannot join the online class.} = « Il n'y a pas de signal dans le village, donc je ne peux pas suivre le cours en ligne. »

\eng{I keep getting disconnected during video calls.} = « Je suis sans cesse déconnecté pendant les appels vidéo. »
\end{exemplebox}

\courssub{Champ 4 — Les difficultés des utilisateurs (user difficulties)}

Tous les problèmes ne viennent pas des machines : parfois, c'est l'utilisateur qui est en difficulté.

\begin{center}
\begin{tabularx}{\linewidth}{|X|X|}\hline
\rowcolor{popblueL} \textbf{English} & \textbf{Français} \\ \hline
to forget a password & oublier un mot de passe \\ \hline
to be locked out & être bloqué, ne plus pouvoir accéder \\ \hline
to be computer-illiterate & être illettré en informatique \\ \hline
to struggle with & avoir du mal avec, peiner sur \\ \hline
to get confused & s'embrouiller, être dérouté \\ \hline
a lack of training & un manque de formation \\ \hline
\end{tabularx}
\end{center}

\eng{She forgot her password and got locked out.} = « Elle a oublié son mot de passe et s'est retrouvée bloquée. »

\eng{Many teachers struggle with new educational software because of a lack of training.} = « Beaucoup d'enseignants peinent avec les nouveaux logiciels éducatifs à cause d'un manque de formation. »

\begin{attention}[Attention aux calques du français]
Ne traduisez pas « je suis confus » par \eng{I am confused} si vous voulez dire « je suis gêné, embarrassé » : \eng{confused} signifie « dérouté, qui ne comprend pas ». Pour « embarrassé », dites \eng{embarrassed}. De même, \eng{to forget a password} : l'ordre des mots est \eng{forget + nom}, jamais \eng{forget of my password}.
\end{attention}

\courssub{Champ 5 — Conséquences et émotions}

Un problème informatique a toujours des conséquences pratiques et émotionnelles.

\begin{center}
\begin{tabularx}{\linewidth}{|X|X|}\hline
\rowcolor{popblueL} \textbf{English} & \textbf{Français} \\ \hline
to lose data & perdre des données \\ \hline
a delay & un retard \\ \hline
frustration (n.) / frustrated (adj.) & la frustration / frustré \\ \hline
helpless (adj.) & impuissant, désemparé \\ \hline
stressful (adj.) & stressant \\ \hline
time-consuming (adj.) & chronophage \\ \hline
to waste time & perdre son temps \\ \hline
\end{tabularx}
\end{center}

\eng{When the computer crashed, I lost all my data and felt completely helpless.} = « Quand l'ordinateur a planté, j'ai perdu toutes mes données et je me suis senti complètement désemparé. »

Remarquez la distinction adjectivale : \eng{stressful} qualifie la \textbf{situation} (\eng{a stressful day}), tandis que \eng{stressed} qualifie la \textbf{personne} (\eng{I feel stressed}). Même logique pour \eng{frustrating} (la situation) et \eng{frustrated} (la personne).

\courssub{Champ 6 — Les solutions}

À chaque problème, sa solution : c'est le champ lexical de la réparation.

\begin{definition}[Troubleshooting]
\cle{Troubleshooting} is the process of finding and fixing problems in a system. En français : le dépannage. Le verbe correspondant est \eng{to troubleshoot} : \eng{The technician is troubleshooting the network.} (« Le technicien dépanne le réseau. »).
\end{definition}

\begin{center}
\begin{tabularx}{\linewidth}{|X|X|}\hline
\rowcolor{popblueL} \textbf{English} & \textbf{Français} \\ \hline
to troubleshoot & dépanner \\ \hline
to fix / to repair & réparer \\ \hline
to restart / to reboot & redémarrer \\ \hline
to call a technician & appeler un technicien \\ \hline
to back up data & sauvegarder des données \\ \hline
technical support / helpline & l'assistance technique \\ \hline
to reset a password & réinitialiser un mot de passe \\ \hline
\end{tabularx}
\end{center}

\eng{Always back up your files.} = « Sauvegardez toujours vos fichiers. »

\eng{Have you tried rebooting the router?} = « As-tu essayé de redémarrer le routeur ? »

La figure suivante relie les problèmes les plus fréquents à leurs solutions naturelles.

\begin{popfigure}
\begin{center}
\begin{tikzpicture}[
  prob/.style={rectangle, rounded corners, draw=popink, fill=poppinkL, align=center, minimum width=4.2cm, minimum height=0.9cm, font=\small},
  sol/.style={rectangle, rounded corners, draw=popgreen, fill=popgreenL, align=center, minimum width=4.2cm, minimum height=0.9cm, font=\small},
  arr/.style={-{Stealth[length=3mm]}, thick, popblue}
]
\node[font=\bfseries\small, popink] at (-3,2.4) {Problem};
\node[font=\bfseries\small, popgreen!60!black] at (3,2.4) {Solution};
\node[prob] (p1) at (-3,1.5) {the computer crashes};
\node[sol] (s1) at (3,1.5) {to reboot / restart it};
\node[prob] (p2) at (-3,0.3) {a virus};
\node[sol] (s2) at (3,0.3) {to install an antivirus};
\node[prob] (p3) at (-3,-0.9) {a forgotten password};
\node[sol] (s3) at (3,-0.9) {to reset the password};
\node[prob] (p4) at (-3,-2.1) {risk of losing data};
\node[sol] (s4) at (3,-2.1) {to back up the files};
\draw[arr] (p1) -- (s1);
\draw[arr] (p2) -- (s2);
\draw[arr] (p3) -- (s3);
\draw[arr] (p4) -- (s4);
\end{tikzpicture}
\end{center}
\legende{Du problème à la solution : les paires lexicales essentielles.}
\end{popfigure}

\begin{aretenir}[L'essentiel des six champs lexicaux]
\begin{itemize}
\item \textbf{Hardware} : \eng{breakdown}, \eng{failure}, \eng{fault}, \eng{out of order}, \eng{dead battery}, \eng{broken screen}.
\item \textbf{Software} : \eng{bug}, \eng{virus}, \eng{crash}, \eng{freeze}, \eng{error message}, \eng{corrupted file}, \eng{to install/uninstall}, \eng{to update}.
\item \textbf{Network} : \eng{slow connection}, \eng{no signal}, \eng{network failure}, \eng{to disconnect}, \eng{offline}, \eng{bandwidth}, \eng{hotspot}.
\item \textbf{Users} : \eng{to forget a password}, \eng{to be locked out}, \eng{computer-illiterate}, \eng{to struggle with}, \eng{lack of training}.
\item \textbf{Consequences} : \eng{to lose data}, \eng{delay}, \eng{frustration}, \eng{helpless}, \eng{stressful}, \eng{time-consuming}, \eng{to waste time}.
\item \textbf{Solutions} : \eng{to troubleshoot}, \eng{to fix}, \eng{to reboot}, \eng{to call a technician}, \eng{to back up}, \eng{helpline}.
\end{itemize}
\end{aretenir}

\begin{exoresolu}[Complétez avec le mot du champ lexical approprié]
Énoncé. Complétez chaque phrase avec un mot ou une expression étudié(e) : \emph{out of order — locked out — back up — crash — bandwidth — troubleshoot}.

1. The photocopier is \ldots\ again; we cannot make any copies.

2. I typed the wrong password three times and got \ldots\ of my e-mail account.

3. Our connection is slow because the \ldots\ is too limited for twenty students.

4. The technician spent two hours trying to \ldots\ the server.

5. If you do not \ldots\ your documents, you may lose everything.

6. My laptop tends to \ldots\ whenever I open too many programs.
\tcblower
Solution détaillée.

1. \eng{out of order} — il s'agit d'une machine (la photocopieuse) en panne.

2. \eng{locked out} — après trois mots de passe faux, l'utilisateur est bloqué.

3. \eng{bandwidth} — la bande passante limitée explique la lenteur de la connexion.

4. \eng{troubleshoot} — le technicien cherche et répare la panne : il dépanne.

5. \eng{back up} — sauvegarder ses documents pour ne pas les perdre.

6. \eng{crash} — l'ordinateur « plante » quand trop de programmes sont ouverts.
\end{exoresolu}

\begin{exercice}[À vous de jouer]
Traduisez en anglais : 1. « L'imprimante est en panne depuis hier. » 2. « Elle a oublié son mot de passe et s'est retrouvée bloquée. » 3. « Sauvegardez toujours vos fichiers avant de mettre à jour le logiciel. » 4. « Le réseau est en panne, nous sommes hors ligne. » Puis rédigez trois phrases personnelles décrivant un problème informatique que vous avez rencontré, en utilisant au moins trois champs lexicaux différents.
\end{exercice}

\begin{corrige}[Exercice — À vous de jouer]
1. \eng{The printer has been out of order since yesterday.}

2. \eng{She forgot her password and got locked out.}

3. \eng{Always back up your files before updating the software.}

4. \eng{The network is down; we are offline.}

Phrases personnelles : réponses libres. Exemple attendu : \eng{Last week, my computer froze during an online test. I felt frustrated because I wasted a lot of time, but the technician fixed it quickly.}
\end{corrige}

Ces six champs lexicaux forment désormais votre boîte à outils. Dans la section suivante, nous verrons comment ces mots se combinent naturellement entre eux grâce aux \cle{collocations} et aux expressions idiomatiques.

\coursec{Collocations et expressions idiomatiques}

Dans la section précédente, nous avons constitué les champs lexicaux des difficultés liées aux TIC : le matériel, les pannes, les logiciels, les connexions. Mais connaître un mot isolé ne suffit pas pour parler un anglais naturel. Un anglophone ne dit jamais \eng{the internet is slow today} sans penser automatiquement \eng{slow connection}, et il ne dit jamais \eng{the system is broken} pour un serveur en panne, mais \eng{the system is down}. Ces associations obligées entre mots s'appellent des \cle{collocations} ; c'est elles qui font la différence entre un anglais correct et un anglais naturel. Cette section vous apprend à associer les mots comme le font les anglophones.

\courssub{Observation : qu'est-ce qu'une collocation ?}

\begin{textebox}[At the computer lab — a short dialogue]
\textbf{Mrs Etame:} What's wrong with your computer, John? The online test starts in five minutes!\\
\textbf{John:} I don't know, Madam. It \textbf{keeps freezing}. I think I \textbf{lost my data} when the system \textbf{crashed} this morning.\\
\textbf{Mrs Etame:} Did you save your file before?\\
\textbf{John:} Yes, I \textbf{backed it up} on a flash drive, but now the computer \textbf{won't start} at all.\\
\textbf{Mrs Etame:} Let me have a look. First, \textbf{reset your password} and try to log in again... No, the screen is black. There must be a \textbf{power failure} in this block.\\
\textbf{John:} Oh no! I'm \textbf{at my wits' end} with this machine!\\
\textbf{Mrs Etame:} Calm down. It's probably just a \textbf{technical hitch}. The technician will \textbf{fix the problem} quickly.\\
\textbf{John:} I hope so. Last week the whole \textbf{network was down} for two days because of a \textbf{major breakdown}.\\
\textbf{Mrs Etame:} Well, look on the bright side: if the test is postponed, you will have more time to revise. It could be a \textbf{blessing in disguise}!
\end{textebox}

\textbf{Glossaire :} \emph{flash drive} (n.) clé USB ; \emph{to log in} (v.) se connecter ; \emph{block} (n.) bâtiment ; \emph{to look on the bright side} (exp.) voir le bon côté des choses ; \emph{to postpone} (v.) reporter.

Observez les groupes de mots en gras du dialogue : \eng{keeps freezing}, \eng{lost my data}, \eng{backed it up}, \eng{won't start}, \eng{reset your password}, \eng{power failure}, \eng{technical hitch}, \eng{fix the problem}, \eng{major breakdown}, \eng{network was down}, \eng{blessing in disguise}. Aucun de ces groupes ne peut être modifié librement : on ne dit pas \eng{make a backup of data} dans ce sens, ni \eng{energetic failure}, ni \eng{the network was broken}. La collocation est un mariage de mots : elle s'apprend comme un tout.

\begin{definition}[Collocation]
Une \cle{collocation} est une association habituelle et quasi obligée de deux mots ou plus (\eng{slow connection}, \eng{to lose data}, \eng{major breakdown}). Le français fonctionne différemment : on ne peut donc pas traduire mot à mot. On distingue les collocations \textbf{verbe + nom} (\eng{to back up a file}), les collocations \textbf{adjectif + nom} (\eng{faulty equipment}) et les \textbf{expressions figées} (\eng{the system is down}).
\end{definition}

\courssub{Collocations verbe + nom}

Voici les verbes qui « vont avec » les noms du domaine informatique. Apprenez chaque couple comme un seul mot.

\begin{center}
\begin{tabularx}{\linewidth}{|X|X|X|}\hline
\rowcolor{popblueL} Collocation & Traduction & Exemple en contexte \\ \hline
to crash (a system) & faire planter (un système) & A virus crashed the whole system. / The system crashed. \\ \hline
to fix a problem & régler un problème & The technician fixed the problem in ten minutes. \\ \hline
to lose data & perdre des données & We lost all our data during the power cut. \\ \hline
to save a file & enregistrer un fichier & Save your file before you close the program. \\ \hline
to back up a file & sauvegarder un fichier (copie de sécurité) & Always back up your documents on a flash drive. \\ \hline
to reset a password & réinitialiser un mot de passe & I forgot my password, so I had to reset it. \\ \hline
to download a document & télécharger un document (vers son appareil) & She downloaded the exam paper from the school website. \\ \hline
to upload a document & téléverser un document (vers un serveur) & Upload your composition before midnight. \\ \hline
to run a program & exécuter un programme & This old computer cannot run the new program. \\ \hline
\end{tabularx}
\end{center}

\begin{propriete}[Download / upload : le sens du mouvement]
\eng{To download} = faire descendre les données \textbf{vers soi} (du serveur vers l'appareil) : \eng{I downloaded the lesson.} « J'ai téléchargé la leçon. » \eng{To upload} = envoyer les données \textbf{vers le serveur} (de l'appareil vers Internet) : \eng{Students must upload their essays.} « Les élèves doivent téléverser leurs dissertations. » Le français « télécharger » couvre souvent les deux sens ; l'anglais, lui, distingue toujours la direction.
\end{propriete}

\begin{exemplebox}[Exemples traduits]
\eng{It keeps freezing.} « Ça n'arrête pas de se figer. »\\
\eng{There was a technical hitch during the online exam.} « Il y a eu un contretemps technique pendant l'examen en ligne. »\\
\eng{Could you help me fix this bug?} « Pourriez-vous m'aider à corriger ce bogue ? »\\
\eng{The lab assistant ran the program on every computer.} « L'assistant de salle informatique a exécuté le programme sur chaque ordinateur. »
\end{exemplebox}

\courssub{Collocations adjectif + nom}

L'adjectif qui qualifie une difficulté technique n'est pas libre non plus. On dit \eng{slow connection} mais jamais \eng{slow internet} seul dans un anglais soigné ; on dit \eng{major breakdown} et non \eng{big breakdown}.

\begin{center}
\begin{tabularx}{\linewidth}{|X|X|X|}\hline
\rowcolor{popblueL} Collocation & Traduction & Exemple en contexte \\ \hline
slow connection & connexion lente & The connection is too slow to watch the video lesson. \\ \hline
technical hitch & contretemps technique & There was a slight technical hitch during the broadcast. \\ \hline
power failure / power cut & panne de courant & A power failure interrupted the online classes. \\ \hline
major breakdown & panne majeure & The server suffered a major breakdown last night. \\ \hline
faulty equipment & équipement défectueux & Much of the equipment in the lab is faulty. \\ \hline
weak signal & signal faible & I can't call you; the signal is very weak in the village. \\ \hline
\end{tabularx}
\end{center}

\begin{attention}[L'adjectif « big » n'est pas universel]
Le français « gros / grand » se traduit rarement par \eng{big} dans les collocations techniques : on dit \eng{a \textbf{major} breakdown}, \eng{a \textbf{serious} problem}, \eng{a \textbf{heavy} user} (gros consommateur), \eng{a \textbf{strong} signal} (et son contraire \eng{weak signal}). Traduire « grosse panne » par \eng{big breakdown} est un calque typique du francophone.
\end{attention}

\courssub{Expressions figées et idiomatiques}

Certaines expressions sont totalement figées : on ne peut ni changer un mot, ni les traduire littéralement.

\begin{propriete}[« The system is down », pas « broken »]
En anglais, on dit \eng{the system is down} et non \eng{the system is broken} pour un réseau, un serveur ou un site web en panne. \cle{Down} est l'adjectif consacré pour les systèmes informatiques ; son contraire est \eng{up} : \eng{The network is down this morning, but it should be up again by noon.} « Le réseau est en panne ce matin, mais il devrait être rétabli d'ici midi. » De même : \eng{the line is bad} (la ligne est mauvaise, au téléphone), \eng{it keeps freezing} (ça n'arrête pas de se figer, avec \eng{keep + V-ing} pour une répétition agaçante).
\end{propriete}

\begin{center}
\begin{tabularx}{\linewidth}{|X|X|}\hline
\rowcolor{popblueL} Expression figée & Traduction et emploi \\ \hline
the system is down & le système est en panne (réseau, serveur) \\ \hline
the line is bad & la ligne est mauvaise (téléphone) \\ \hline
it keeps freezing & ça n'arrête pas de se figer (écran, ordinateur) \\ \hline
it won't start & il refuse de démarrer (volonté négative : won't) \\ \hline
it doesn't respond & il ne répond pas (appareil bloqué) \\ \hline
to be at a loss & être désemparé, ne pas savoir quoi faire \\ \hline
to be at one's wits' end & être au bout du rouleau \\ \hline
a blessing in disguise & un mal pour un bien \\ \hline
\end{tabularx}
\end{center}

\begin{exemplebox}[Exemples traduits]
\eng{I'm at my wits' end with this laptop.} « Je suis au bout du rouleau avec cet ordinateur portable. »\\
\eng{When the file disappeared, the secretary was at a loss.} « Quand le fichier a disparu, la secrétaire était désemparée. »\\
\eng{Sorry, I can't hear you — the line is very bad.} « Désolé, je ne t'entends pas, la ligne est très mauvaise. »\\
\eng{Missing the online class was a blessing in disguise: the recording was much clearer.} « Rater le cours en ligne a été un mal pour un bien : l'enregistrement était bien plus clair. »
\end{exemplebox}

Notez la forme de \eng{to be at one's wits' end} : le possessif change selon le sujet (\eng{my wits' end}, \eng{her wits' end}) et \eng{wits} prend un apostrophe après le s. Quant à \eng{it won't start}, l'anglais utilise ici \eng{won't} (volonté refusée) pour une machine capricieuse : c'est une personnification très fréquente — \eng{My phone won't charge.} « Mon téléphone refuse de charger. »

\courssub{Demander de l'aide et décrire une panne}

\begin{exemplebox}[Formules utiles à l'oral]
\eng{Could you help me fix the printer?} « Pourriez-vous m'aider à réparer l'imprimante ? »\\
\eng{What's wrong with my computer?} « Qu'est-ce qui ne va pas avec mon ordinateur ? »\\
\eng{It won't start.} « Il ne démarre pas. »\\
\eng{It doesn't respond when I click.} « Il ne répond pas quand je clique. »\\
\eng{The screen has gone blank.} « L'écran est devenu noir. »
\end{exemplebox}

\begin{attention}[Éviter la phrase fourre-tout]
Ne dites pas \eng{the computer has a problem} à tout propos : c'est vague et c'est un calque de « l'ordinateur a un problème ». Préférez les collocations précises : \eng{the computer crashed} (il a planté), \eng{the computer is frozen} (il est figé), \eng{the computer won't start} (il ne démarre pas), \eng{the computer is very slow} (il est très lent). La précision lexicale est un critère d'évaluation à l'examen.
\end{attention}

\courssub{To fix, to repair ou to mend ?}

Ces trois verbes signifient « réparer », mais ils ne s'emploient pas dans les mêmes contextes.

\begin{center}
\begin{tabularx}{\linewidth}{|X|X|X|}\hline
\rowcolor{popblueL} Verbe & Emploi & Exemple \\ \hline
to fix & général ; régler un problème, réparer un appareil ; aussi « fixer » (date, réunion) ; courant en AmE et BrE & Can you fix my laptop? Let's fix a date for the meeting. \\ \hline
to repair & technique, soutenu ; réparation professionnelle (surtout BrE) & The technician repaired the server. \\ \hline
to mend & objets simples de la vie quotidienne (vêtements, chaussures, vélo) ; BrE, vieillit pour les appareils électroniques & She mended her torn uniform. \\ \hline
\end{tabularx}
\end{center}

\begin{popfigure}
\begin{center}
\begin{tikzpicture}
\node[circle, draw=popblue, fill=popblueL, thick, minimum size=2.2cm, align=center, font=\bfseries] (fix) at (0,0) {to fix};
\node[draw=poporange, fill=poporangeL, rounded corners, thick, align=center] (bug) at (-4,2) {a bug};
\node[draw=popgreen, fill=popgreenL, rounded corners, thick, align=center] (comp) at (0,3) {a computer};
\node[draw=poppurple, fill=poppurpleL, rounded corners, thick, align=center] (prob) at (4,2) {a problem};
\node[draw=poppink, fill=poppinkL, rounded corners, thick, align=center] (date) at (-3.2,-2.2) {a date};
\node[draw=popgold, fill=popgoldL, rounded corners, thick, align=center] (meet) at (3.2,-2.2) {a meeting};
\draw[-{Stealth}, thick, popblue] (fix) -- (bug);
\draw[-{Stealth}, thick, popblue] (fix) -- (comp);
\draw[-{Stealth}, thick, popblue] (fix) -- (prob);
\draw[-{Stealth}, thick, popblue] (fix) -- (date);
\draw[-{Stealth}, thick, popblue] (fix) -- (meet);
\end{tikzpicture}
\end{center}
\legende{La polyvalence de « to fix » : réparer un bogue ou un ordinateur, régler un problème, mais aussi fixer une date ou une réunion.}
\end{popfigure}

\begin{exoresolu}[Choisissez la bonne collocation]
Complétez avec la collocation correcte : \emph{down / backed up / reset / weak / fix / at my wits' end / hitch / won't}.\\
1. The school website is \dots\ again; nobody can check the results.\\
2. I \dots\ my project on a flash drive before the power cut.\\
3. I forgot my password, so the administrator had to \dots\ it.\\
4. The signal is too \dots\ here to make a video call.\\
5. Could you help me \dots\ this problem with the printer?\\
6. I'm \dots\ with this old computer!\\
7. There was a small technical \dots\ at the beginning of the online lesson.\\
8. My phone \dots\ charge; I think the battery is dead.
\tcblower
1. \eng{down} — le site est en panne (adjectif consacré pour les systèmes). 2. \eng{backed up} — sauvegardé (copie de sécurité). 3. \eng{reset} — réinitialiser un mot de passe. 4. \eng{weak} — signal faible. 5. \eng{fix} — régler un problème. 6. \eng{at my wits' end} — au bout du rouleau (possessif \eng{my} avec le sujet \eng{I}). 7. \eng{hitch} — contretemps technique. 8. \eng{won't} — le téléphone « refuse » de charger (personnification).
\end{exoresolu}

\begin{methode}[Comment mémoriser les collocations]
\begin{enumerate}
\item N'apprenez jamais un mot seul : apprenez-le avec son partenaire habituel — \eng{slow + connection} (et non « slow + internet » seul), \eng{back up + a file}, \eng{reset + a password}.
\item Apprenez chaque collocation dans une phrase complète et personnelle : \eng{I lost my data during the power cut in Douala.} Une phrase vraie se retient mieux qu'une liste.
\item Classez votre carnet de vocabulaire par verbes : une page « to fix », une page « to lose », avec tous leurs partenaires.
\item Réemployez chaque collocation à l'oral dans la semaine : c'est l'usage qui la fixe durablement.
\end{enumerate}
\end{methode}

\begin{exercice}[À vous de jouer]
1. Traduisez en anglais : « Le réseau est en panne, alors je ne peux pas téléverser ma composition. » / « Je suis désemparé : mon ordinateur n'arrête pas de se figer. »\\
2. Trouvez la collocation qui ne convient pas : (a) \emph{major} breakdown / \emph{big} breakdown ; (b) \emph{weak} signal / \emph{feeble} signal ; (c) \emph{technical} hitch / \emph{technical} problem.\\
3. Rédigez un court dialogue (6 répliques) entre un élève et un technicien, en utilisant au moins cinq collocations de cette section.
\end{exercice}

\begin{corrige}[Exercice — éléments de réponse]
1. \eng{The network is down, so I can't upload my composition.} / \eng{I'm at a loss: my computer keeps freezing.}\\
2. (a) \eng{big breakdown} est incorrect : on dit \eng{major breakdown}. (b) \eng{feeble signal} est incorrect : on dit \eng{weak signal}. (c) les deux existent, mais \eng{technical hitch} désigne un petit contretemps passager.\\
3. Modèle : \eng{— What's wrong with your tablet? — It won't start and the screen is frozen. — Did you back up your files? — Yes, luckily. — Good. I'll fix the problem and reset your password. — Thank you so much!}
\end{corrige}

\coursec{Faux amis et calques du français}

Après avoir exploré les collocations et expressions idiomatiques qui structurent le discours sur les pannes, la lenteur et la sécurité informatique, nous devons à présent nous prémunir contre un danger sournois : l'interférence de la langue maternelle. Le français et l'anglais partagent des milliers de mots d'origine commune (les « vrais amis »), mais certains ont divergé au fil des siècles. Dans le domaine des TIC, où l'anglais est la langue de référence, confondre \eng{deception} et « tromperie » ou \eng{actually} et « actuellement » trahit immédiatement un locuteur non natif et peut changer radicalement le sens d'un message professionnel ou académique. Cette section identifie les pièges lexicaux et syntaxiques les plus fréquents pour un francophone camerounais.

\courssub{Les cinq faux amis majeurs du thème TIC}

Commençons par observer ces mots dans des phrases authentiques tirées de forums d'entraide technique, de mails professionnels et de rapports d'étudiants de l'Université de Buea ou de Yaoundé I.

\begin{textebox}[Extraits authentiques : faux amis en contexte]
\textbf{1.} “The system crash was a major \textbf{deception} for the team who had worked all night.” \\
\textbf{2.} “\textbf{Actually}, the new software update solved the compatibility issue we had last week.” \\
\textbf{3.} “Senior technicians \textbf{assist} junior staff during the network configuration process.” \\
\textbf{4.} “The client \textbf{demanded} a full refund after the data loss incident.” \\
\textbf{5.} “Students often go to the university \textbf{library} to access online journals for their research.” \\
\textbf{6.} “I need to \textbf{ask} the IT manager for a new password.” \\
\textbf{7.} “She \textbf{attends} a coding bootcamp every Saturday in Douala.”
\end{textebox}

\begin{definition}[Faux ami (False friend)]
Un \cle{faux ami} est un mot anglais qui ressemble à s'y méprendre à un mot français (orthographe ou prononciation proches) mais dont le sens a divergé. L'utilisation du sens français dans une phrase anglaise produit une erreur lexicale grave, souvent incompréhensible pour un anglophone.
\end{definition}

Analysons chaque paire de manière rigoureuse.

\begin{propriete}[Faux ami 1 : \eng{deception} $\neq$ tromperie]
\textbf{Forme} : \eng{deception} (nom, non comptable ou comptable). \\
\textbf{Sens réel} : \cle{déception}, \cle{désillusion}, sentiment de tristesse parce que quelque chose n'est pas aussi bon qu'espéré. \\
\textbf{Piège français} : \eng{deception} ne signifie \textbf{jamais} « tromperie », « fraude » ou « arnaque ». \\
\textbf{Mots corrects pour « tromperie »} : \eng{deceit} (n.), \eng{fraud} (n.), \eng{scam} (n., familier/informel), \eng{trickery} (n.). \\
\textbf{Verbe associé} : \eng{to deceive} (tromper, duper) $\rightarrow$ \eng{deceptive} (adj., trompeur). \\
\textbf{Exemples} : \eng{The slow internet speed was a huge deception for the online candidates.} (La lenteur d'internet a été une énorme déception pour les candidats en ligne.) $\neq$ \eng{The email was a fraud/scam.} (L'email était une arnaque.)
\end{propriete}

\begin{propriete}[Faux ami 2 : \eng{actually} $\neq$ actuellement]
\textbf{Forme} : \eng{actually} (adverbe de modalité/opposition). \\
\textbf{Sens réel} : \cle{en fait}, \cle{réellement}, \cle{en réalité}. Il marque une correction, une surprise ou une précision par rapport à ce qui précède. \\
\textbf{Piège français} : \eng{actually} ne signifie \textbf{jamais} « actuellement », « de nos jours », « pour l'instant ». \\
\textbf{Mots corrects pour « actuellement »} : \eng{currently}, \eng{at present}, \eng{at the moment}, \eng{nowadays}. \\
\textbf{Exemples} : \eng{People think the server is down, but actually it is just rebooting.} (Les gens pensent que le serveur est en panne, mais en fait il redémarre simplement.) $\neq$ \eng{Currently, we are upgrading the firewall.} (Actuellement, nous mettons à jour le pare-feu.)
\end{propriete}

\begin{propriete}[Faux ami 3 : \eng{to assist} $\neq$ assister à]
\textbf{Forme} : \eng{to assist} (verbe transitif direct, registre formel/professionnel). \\
\textbf{Sens réel} : \cle{aider}, \cle{prêter assistance à} quelqu'un. \\
\textbf{Piège français} : \eng{to assist} ne signifie \textbf{jamais} « assister à » (être présent à un événement, un cours, une réunion). \\
\textbf{Mot correct pour « assister à »} : \eng{to attend} (verbe transitif direct). \\
\textbf{Distinction cruciale} : \eng{to assist somebody} (aider quelqu'un) vs \eng{to attend something} (assister à quelque chose : cours, meeting, conférence). \\
\textbf{Exemples} : \eng{The helpdesk assists users with password resets.} (Le service d'assistance aide les utilisateurs à réinitialiser leurs mots de passe.) / \eng{All Form Five students must attend the ICT workshop.} (Tous les élèves de Première doivent assister à l'atelier TIC.)
\end{propriete}

\begin{propriete}[Faux ami 4 : \eng{to demand} $\neq$ demander]
\textbf{Forme} : \eng{to demand} (verbe transitif, fort). \\
\textbf{Sens réel} : \cle{exiger}, \cle{réclamer avec autorité}. Cela implique un rapport de force, une absence de négociation. \\
\textbf{Piège français} : \eng{to demand} ne s'emploie \textbf{pas} pour le sens neutre de « demander » (une information, un service, poliment). \\
\textbf{Mots corrects pour « demander »} : \eng{to ask (for)}, \eng{to request} (formel), \eng{to inquire about} (se renseigner sur). \\
\textbf{Exemples} : \eng{The union demanded better cybersecurity measures.} (Le syndicat a exigé de meilleures mesures de cybersécurité.) / \eng{I asked the technician for a quick diagnosis.} (J'ai demandé au technicien un diagnostic rapide.)
\end{propriete}

\begin{propriete}[Faux ami 5 : \eng{a library} $\neq$ une librairie]
\textbf{Forme} : \eng{library} (nom comptable). \\
\textbf{Sens réel} : \cle{bibliothèque} (lieu de prêt/consultation de livres, centre de ressources, \eng{digital library}). \\
\textbf{Piège français} : \eng{library} ne désigne \textbf{jamais} un commerce vendant des livres. \\
\textbf{Mot correct pour « librairie »} : \eng{bookshop} (UK), \eng{bookstore} (US). \\
\textbf{Exemples} : \eng{The school library has a new e-learning section.} (La bibliothèque de l'école a une nouvelle section d'e-learning.) / \eng{I bought this programming manual at the bookshop near the market.} (J'ai acheté ce manuel de programmation à la librairie près du marché.)
\end{propriete}

\begin{center}
\begin{tabularx}{\linewidth}{|X|X|X|}
\hline
\rowcolor{popblueL} \textbf{Mot anglais} & \textbf{Sens réel (Français)} & \textbf{Piège français (faux sens)} \\ \hline
\eng{deception} & Déception, désillusion & $\neq$ Tromperie, fraude (\eng{deceit, fraud}) \\ \hline
\eng{actually} & En fait, réellement & $\neq$ Actuellement (\eng{currently, at present}) \\ \hline
\eng{to assist} & Aider (qn) & $\neq$ Assister à (qch) (\eng{to attend}) \\ \hline
\eng{to demand} & Exiger & $\neq$ Demander (\eng{to ask, to request}) \\ \hline
\eng{library} & Bibliothèque & $\neq$ Librairie (\eng{bookshop/bookstore}) \\ \hline
\end{tabularx}
\end{center}
\textit{Les 5 faux amis critiques du thème TIC.}

\begin{attention}[Erreur classique au Bac : \eng{actually} vs \eng{currently}]
Ne dites \textbf{jamais} : \eng{*“I am actually in Lower Sixth.”} pour dire « Je suis actuellement en Première (Lower Sixth) ».
\eng{Actually} signifie « en fait / en réalité ». La phrase anglaise veut dire : « En fait, je suis en Première (contrairement à ce que tu penses) ».
\par \textbf{Correct} : \eng{“I am \textbf{currently} in Lower Sixth.”} ou \eng{“I am in Lower Sixth \textbf{at the moment}.”}
\par Astuce mnémotechnique : \eng{Actu\underline{ally}} $\rightarrow$ \underline{Real}ly (Réallement). \eng{Curr\underline{ently}} $\rightarrow$ \underline{Current} time (Temps présent).
\end{attention}

\begin{exemplebox}[Exemples traduits : discrimination des faux amis]
\begin{enumerate}
\item \eng{The frequent power cuts were a real \textbf{deception} for the data entry team.} \\
« Les coupures de courant fréquentes ont été une vraie \textbf{déception} pour l'équipe de saisie. »
\item \eng{Many parents think ICT is just games, but \textbf{actually} it teaches coding.} \\
« Beaucoup de parents pensent que les TIC sont juste des jeux, mais \textbf{en réalité} ça apprend le codage. »
\item \eng{The IT technician will \textbf{assist} you with the printer installation.} \\
« Le technicien informatique va vous \textbf{aider} pour l'installation de l'imprimante. »
\item \eng{All trainees must \textbf{attend} the cybersecurity awareness session.} \\
« Tous les stagiaires doivent \textbf{assister à} la session de sensibilisation à la cybersécurité. »
\item \eng{The manager \textbf{demanded} an immediate explanation for the server crash.} \\
« Le responsable a \textbf{exigé} une explication immédiate pour la panne du serveur. »
\item \eng{Please \textbf{ask} the receptionist for the Wi-Fi password.} \\
« Veuillez \textbf{demander} à la réceptionniste le mot de passe Wi-Fi. »
\item \eng{I borrowed a textbook on Python from the university \textbf{library}.} \\
« J'ai emprunté un manuel sur Python à la \textbf{bibliothèque} universitaire. »
\item \eng{She bought a dictionary at the \textbf{bookshop} in Mfoundi market.} \\
« Elle a acheté un dictionnaire à la \textbf{librairie} du marché Mfoundi. »
\end{enumerate}
\end{exemplebox}

\courssub{Calques syntaxiques et erreurs de structure}

Au-delà du lexique, la structure de la phrase est souvent calquée sur le français. Deux erreurs reviennent systématiquement dans les productions d'élèves camerounais lorsqu'ils décrivent des problèmes techniques.

\begin{propriete}[Structure : \eng{to have a problem with} vs *\eng{to have a problem for}]
\textbf{Règle} : En anglais, la préposition qui suit \eng{problem} (ou \eng{difficulty, trouble, issue}) pour introduire la source du problème est \textbf{\eng{with}}.
\par \textbf{Structure} : \eng{Subject + have/has + a problem/difficulty/issue + with + Noun/Gerund}.
\par \textbf{Calque fautif} : *\eng{“My computer has a problem for me.”} (Traduction littérale de « Mon ordinateur a un problème pour moi »).
\par \textbf{Correct} : \eng{“I \textbf{have a problem \textbf{with}} my computer.”} ou \eng{“My computer \textbf{has a problem}.”} (sans complément d'agent).
\par \textbf{Variantes} : \eng{to experience a problem with}, \eng{to encounter an issue with}, \eng{to face difficulties with}.
\end{propriete}

\begin{propriete}[Adjectif vs Adverbe : \eng{slow} vs \eng{slowly}]
\textbf{Règle} : \eng{Slow} est un \textbf{adjectif} (qualifie un nom). \eng{Slowly} est un \textbf{adverbe} (qualifie un verbe, un adjectif ou un autre adverbe).
\par Après un verbe d'état (\eng{be, seem, become, look, feel, run, work}), on utilise l'\textbf{adjectif}.
\par \textbf{Calque fautif} : *\eng{“The connection is \textbf{slowly}.”} (Calque sur « La connexion est lentement » $\leftarrow$ confusion avec l'adverbe français).
\par \textbf{Correct} : \eng{“The connection \textbf{is slow}.”} (La connexion est lente).
\par \textbf{Adverbe correct} : \eng{“The file downloads \textbf{slowly}.”} (Le fichier se télécharge lentement — modifie le verbe \eng{downloads}).
\end{propriete}

\begin{exemplebox}[Correction de calques syntaxiques]
\begin{center}
\begin{tabularx}{\linewidth}{|X|X|}\hline
\rowcolor{popblueL} \textbf{Phrase fautive (Calque FR)} & \textbf{Phrase correcte (English)} \\ \hline
*\eng{My phone has a problem for me.} & \eng{I have a problem \textbf{with} my phone.} \\ \hline
*\eng{The network is slowly today.} & \eng{The network \textbf{is slow} today.} \\ \hline
*\eng{We have difficulties to connect.} & \eng{We have difficulties \textbf{connecting} / \textbf{to connect}.} \\ \hline
*\eng{He demanded me a new cable.} & \eng{He \textbf{asked me for} a new cable.} (\eng{demand} ne prend pas COI) \\ \hline
*\eng{She assists the meeting.} & \eng{She \textbf{attends} the meeting.} \\ \hline
\end{tabularx}
\end{center}
\end{exemplebox}

\courssub{L'expression figée : \eng{digital divide}}

Certaines associations de mots (collocations) sont des « blocs » indissociables. Tenter de les traduire mot à mot depuis le français produit un anglais approximatif.

\begin{propriete}[Bloc lexical : \eng{digital divide}]
\textbf{Forme} : \eng{digital divide} (nom composé, souvent avec \eng{the}).
\textbf{Sens} : \cle{fracture numérique} (inégalité d'accès aux TIC entre groupes sociaux, zones géographiques, générations).
\textbf{Piège} : Ne pas dire *\eng{digital division}, *\eng{numeric fracture}, *\eng{digital gap} (bien que \eng{gap} existe, \eng{divide} est le terme technique consacré par l'ONU, l'UIT et les rapports officiels camerounais comme le \eng{National ICT Strategic Plan}).
\textbf{Collocations fréquentes} : \eng{bridge the digital divide} (combler la fracture numérique), \eng{widen the digital divide} (creuser la fracture), \eng{rural/urban digital divide}, \eng{gender digital divide}.
\textbf{Exemples} : \eng{The government aims to \textbf{bridge the digital divide} in rural areas like the Far North.} (Le gouvernement vise à combler la fracture numérique dans les zones rurales comme l'Extrême-Nord.) / \eng{Lack of electricity \textbf{widens the digital divide}.} (Le manque d'électricité creuse la fracture numérique.)
\end{propriete}

\begin{savaistu}[Le faux ami au Baccalauréat]
À l'épreuve écrite du Baccalauréat (Compréhension et Expression), l'utilisation d'un faux ami est sanctionnée comme une \textbf{erreur lexicale majeure} (souvent 0,5 à 1 point en moins par occurrence selon la grille). De même, en Compréhension de l'écrit, ne pas reconnaître \eng{actually} = « en fait » ou \eng{library} = « bibliothèque » entraîne une méprise du sens global du texte. Apprendre ces 5 paires par cœur est un « retour sur investissement » garanti pour votre note.
\end{savaistu}

\begin{aretenir}[Synthèse : Les 5 faux amis + 2 pièges syntaxiques à maîtriser]
\begin{itemize}
\item \eng{deception} = \textbf{déception} (pas tromperie $\rightarrow$ \eng{fraud/scam}).
\item \eng{actually} = \textbf{en fait} (pas actuellement $\rightarrow$ \eng{currently/at present}).
\item \eng{to assist} = \textbf{aider} (pas assister à $\rightarrow$ \eng{to attend}).
\item \eng{to demand} = \textbf{exiger} (pas demander $\rightarrow$ \eng{to ask/request}).
\item \eng{library} = \textbf{bibliothèque} (pas librairie $\rightarrow$ \eng{bookshop}).
\item \textbf{Problème} : \eng{have a problem \textbf{with}} (pas *\eng{for}).
\item \textbf{Lent} : \eng{is \textbf{slow}} (adj.) vs \eng{runs \textbf{slowly}} (adv.).
\end{itemize}
\end{aretenir}

\courssub{Entraînement guidé : Chasse aux erreurs}

\begin{exoresolu}[Correction d'un paragraphe truffé de faux amis et calques]
\textbf{Consigne} : Le paragraphe ci-dessous a été écrit par un élève de Terminale A. Il contient \textbf{6 erreurs} (faux amis, calques syntaxiques, confusion adjectif/adverbe). Identifiez-les, soulignez-les et réécrivez le paragraphe correctement.

\textbf{Texte de l'élève} :
\eng{“Actually, I am currently in Upper Sixth at Government Bilingual High School Yaoundé. Yesterday, I went to the library to buy a book on networking, but the shop was closed. My computer has a problem for me since Monday: the internet connection is slowly. I demanded my father to assist the repair shop, but he was busy. It was a big deception.”}

\tcblower

\textbf{Solution détaillée et commentée} :

\begin{enumerate}
\item \eng{\textbf{Actually} I am currently...} $\rightarrow$ \eng{\textbf{Currently} / \textbf{At present}, I am...} \\
\textit{Explication} : L'élève veut dire « Actuellement ». \eng{Actually} = « En fait ». Pas de contradiction ici, juste une indication temporelle.
\item \eng{...went to the \textbf{library} to buy a book...} $\rightarrow$ \eng{...went to the \textbf{bookshop} / \textbf{bookstore} to buy a book...} \\
\textit{Explication} : On achète dans une \eng{bookshop}, on emprunte/consulte à la \eng{library}.
\item \eng{My computer \textbf{has a problem for me} since Monday} $\rightarrow$ \eng{I \textbf{have had a problem \textbf{with}} my computer \textbf{since}} Monday. \\
\textit{Explication} : Double correction. 1) Calque syntaxique : \eng{have a problem WITH}. 2) Temps verbal : \eng{since Monday} (point de départ dans le passé) appelle le \textbf{Present Perfect} (\eng{have had}), pas le \eng{Present Simple}. \eng{For me} est inutile.
\item \eng{...connection is \textbf{slowly}.} $\rightarrow$ \eng{...connection \textbf{is slow}.} \\
\textit{Explication} : Après verbe d'état \eng{is}, adjectif \eng{slow}. \eng{Slowly} est adverbe.
\item \eng{I \textbf{demanded} my father to \textbf{assist} the repair shop...} $\rightarrow$ \eng{I \textbf{asked} my father to \textbf{take} / \textbf{bring} the computer to the repair shop...} \\
\textit{Explication} : Double faux ami. 1) \eng{Demanded} = « exigé » (trop fort, et \eng{demand} ne prend pas de COI direct *\eng{demanded my father} ; il faut \eng{demanded \textbf{that} my father...} ou \eng{asked my father}). 2) \eng{Assist} = « aider ». Le père n'aide pas le magasin, il y \textbf{apporte} l'ordinateur. Verbe : \eng{take/bring}.
\item \eng{It was a big \textbf{deception}.} $\rightarrow$ \eng{It was a big \textbf{disappointment} / \textbf{let-down}.} \\
\textit{Explication} : Nuance. \eng{Deception} = sentiment d'être trompé (vol, arnaque). Ici, c'est de la \textbf{déception} face à une situation : \eng{disappointment} s'accorde mieux avec \eng{big/huge} pour une situation technique.
\end{enumerate}

\textbf{Version corrigée} :
\eng{“\textbf{Currently}, I am in Upper Sixth at GBHS Yaoundé. Yesterday, I went to the \textbf{bookshop} to buy a book on networking, but it was closed. \textbf{I have had a problem \textbf{with}} my computer \textbf{since} Monday: the internet connection \textbf{is slow}. \textbf{I asked} my father to \textbf{take the computer to} the repair shop, but he was busy. It was a big \textbf{disappointment}.”}
\end{exoresolu}

\begin{exercice}[Entraînement : Traduction ciblée (Français $\rightarrow$ Anglais)]
Traduisez les phrases suivantes en évitant tous les pièges vus dans cette section (faux amis, calques, adjectif/adverbe).
\begin{enumerate}
\item Actuellement, la connexion Wi-Fi est lente dans l'amphithéâtre.
\item Cette panne de serveur a été une vraie déception pour l'équipe projet.
\item Le technicien a aidé l'enseignant à configurer le vidéoprojecteur.
\item Les élèves doivent assister à la formation sur la cybersécurité la semaine prochaine.
\item J'ai demandé au proviseur d'acheter de nouveaux ordinateurs pour la bibliothèque.
\item La fracture numérique s'agrandit entre les zones urbaines et rurales.
\item Mon téléphone a un problème : il se charge très lentement.
\item En réalité, le logiciel est gratuit pour les étudiants.
\end{enumerate}
\end{exercice}

\begin{corrige}[Exercice : Traduction ciblée]
\begin{enumerate}
\item \eng{\textbf{Currently} / \textbf{At the moment}, the Wi-Fi connection \textbf{is slow} in the lecture hall.} (\textit{Pièges : currently $\neq$ actually ; slow adj $\neq$ slowly}).
\item \eng{This server crash \textbf{was a real disappointment} for the project team.} (\textit{Piège : deception = déception mais disappointment mieux pour une panne ; crash = panne brutale}).
\item \eng{The technician \textbf{assisted} the teacher \textbf{in configuring} / \textbf{to configure} the video projector.} (\textit{Assist = aider ; registre formel. Préposition : assist sb in doing sth}).
\item \eng{Students must \textbf{attend} the cybersecurity training next week.} (\textit{Piège : attend $\neq$ assist}).
\item \eng{I \textbf{asked} the Principal / \textbf{Headmaster} \textbf{for} new computers for the \textbf{library}.} (\textit{Pièges : ask for $\neq$ demand ; library = bibliothèque ; bookshop = librairie}).
\item \eng{The \textbf{digital divide} \textbf{is widening} between urban and rural areas.} (\textit{Bloc lexical : digital divide ; widen = agrandir/creuser}).
\item \eng{I \textbf{have a problem \textbf{with}} my phone: it charges \textbf{very slowly}.} (\textit{Pièges : have a problem with ; slowly = adverbe modifie charges}).
\item \eng{\textbf{Actually}, the software \textbf{is free for} students.} (\textit{Piège : actually = en réalité ; free for}).
\end{enumerate}
\end{corrige}

\coursec{Familles de mots, prononciation et orthographe}

Dans la section précédente, nous avons débusqué les \cle{faux amis} et les calques du français qui guettent l'élève francophone dès qu'il parle des difficultés informatiques. Nous savons désormais qu'on ne dit pas \eng{*the computer is planté} ni \eng{*I have a problem of connexion}. Mais connaître un mot isolé ne suffit pas : en anglais comme en français, les mots vivent en \cle{familles}. Celui qui connaît \eng{connect} peut deviner \eng{connection}, \eng{connectivity}, \eng{disconnected} — à condition de maîtriser les mécanismes de formation des mots, leur prononciation et leur orthographe. C'est l'objet de cette section : apprendre à \emph{fabriquer}, \emph{prononcer} et \emph{écrire} correctement tout le vocabulaire des pannes et des difficultés liées aux TIC.

\courssub{Observer avant de mémoriser : un texte pour découvrir les familles de mots}

\begin{textebox}[A frustrating Monday at the cybercafé]
Last Monday, I went to a cybercafé in Buea to submit my online application before the deadline. Everything went wrong. First, the connection failed three times, so I had to disconnect and reconnect the router. Then the whole system broke down: the screen froze, and the technician told me the server had crashed during the night. “The breakdown is serious,” he said. “We need to restart, reset and reinstall the software.”

While waiting, I started troubleshooting with my phone, but the connectivity was so poor that even WhatsApp refused to work. One computer was completely broken; another one was unavailable because a user had forgotten to log out. The software itself was not user-friendly at all, and an old bug kept interrupting the registration process. “This program is useless,” complained a student next to me. “It is full of malfunctions.”

After two hours, the technician finally disabled the faulty program, updated the system and fixed the bug. “Everything is up-to-date now,” he announced with a smile. I typed my log-in and password, and the page opened at once. What a relief! I submitted my application just in time. That day, I learned an important lesson: technology is useful, but only when it works — and when the user knows what to do when it fails.
\end{textebox}

\textbf{Glossaire}

\begin{center}
\begin{tabularx}{\linewidth}{|l|l|X|}
\hline
\rowcolor{popblueL} Mot anglais & Nature & Traduction française \\
\hline
connection & nom & connexion \\
\hline
to disconnect / reconnect & verbe & déconnecter / reconnecter \\
\hline
to break down & verbe à particule & tomber en panne \\
\hline
breakdown & nom & panne \\
\hline
to freeze (froze, frozen) & verbe irrégulier & se figer, se bloquer (écran) \\
\hline
to crash & verbe & planter (système) \\
\hline
to restart / reset / reinstall & verbes & redémarrer / réinitialiser / réinstaller \\
\hline
connectivity & nom & connectivité (qualité de la connexion) \\
\hline
unavailable & adjectif & indisponible \\
\hline
user-friendly & adjectif composé & convivial, facile d'usage \\
\hline
bug & nom & bogue, bug \\
\hline
useless & adjectif & inutile \\
\hline
malfunction & nom & dysfonctionnement \\
\hline
to disable & verbe & désactiver \\
\hline
up-to-date & adjectif composé & à jour \\
\hline
log-in & nom composé & identifiant de connexion \\
\hline
\end{tabularx}
\end{center}

\textbf{Questions de compréhension}

\begin{enumerate}
\item How many times did the connection fail?
\item What did the technician do to repair the system?
\item Why was one computer unavailable?
\item What did the narrator learn at the end of the day?
\end{enumerate}

\begin{corrige}[Questions de compréhension]
\begin{enumerate}
\item The connection failed \textbf{three times}.
\item He \textbf{disabled} the faulty program, \textbf{updated} the system and \textbf{fixed} the bug (after restarting, resetting and reinstalling the software).
\item One computer was unavailable because \textbf{a user had forgotten to log out}.
\item He learned that technology is useful \textbf{only when it works} — and when the user knows what to do when it fails.
\end{enumerate}
\end{corrige}

\courssub{Les grandes familles de mots du lexique des TIC}

\begin{definition}[Famille de mots]
Une \cle{famille de mots} (\eng{word family}) est un ensemble de mots construits sur la même \cle{racine} (\eng{root}) à l'aide de \cle{préfixes} (placés devant : \eng{dis-}, \eng{re-}) et de \cle{suffixes} (placés derrière : \eng{-ion}, \eng{-ful}, \eng{-less}). Chaque suffixe ou préfixe modifie la nature grammaticale ou le sens du mot de départ.
\end{definition}

\textbf{1. La famille de \eng{connect}}

\begin{center}
\begin{tabularx}{\linewidth}{|l|l|X|}
\hline
\rowcolor{popblueL} Mot & Nature et formation & Traduction et exemple \\
\hline
to connect & verbe (racine) & connecter — \eng{Connect the cable to the router.} « Connecte le câble au routeur. » \\
\hline
connection & nom (connect + -ion) & connexion — \eng{The connection is very slow today.} « La connexion est très lente aujourd'hui. » \\
\hline
connectivity & nom (connect + -ivity) & connectivité — \eng{Rural areas often suffer from poor connectivity.} « Les zones rurales souffrent souvent d'une mauvaise connectivité. » \\
\hline
connected & adjectif (connect + -ed) & connecté — \eng{All the computers are connected to the network.} « Tous les ordinateurs sont connectés au réseau. » \\
\hline
disconnected & adjectif (dis- + connected) & déconnecté — \eng{My laptop is disconnected from the Wi-Fi.} « Mon portable est déconnecté du Wi-Fi. » \\
\hline
to disconnect & verbe (dis- + connect) & déconnecter — \eng{Disconnect the printer before the storm.} « Déconnecte l'imprimante avant l'orage. » \\
\hline
\end{tabularx}
\end{center}

\begin{popfigure}
\begin{tikzpicture}[
  grow=right,
  level 1/.style={sibling distance=2.1cm, level distance=3.6cm},
  every node/.style={draw, rounded corners, align=center, font=\small, inner sep=4pt},
  edge from parent/.style={draw=popmuted, thick}
]
\node[fill=popblue, text=white, font=\small\bfseries] {connect}
  child { node[fill=popgreenL, align=center]{-ed \\ connected} }
  child { node[fill=popgoldL, align=center]{dis- \\ disconnect} }
  child { node[fill=poppinkL, align=center]{-ivity \\ connectivity} }
  child { node[fill=poporangeL, align=center]{-ion \\ connection} };
\end{tikzpicture}
\legende{Arbre lexical : la racine \eng{connect} et ses dérivés formés par préfixes et suffixes.}
\end{popfigure}

\textbf{2. Les familles de \eng{fail} et de \eng{break}}

\begin{center}
\begin{tabularx}{\linewidth}{|l|l|X|}
\hline
\rowcolor{popblueL} Mot & Nature & Traduction et exemple \\
\hline
to fail & verbe & échouer, tomber en panne — \eng{The connection failed during the exam.} « La connexion a échoué pendant l'examen. » \\
\hline
failure & nom (fail + -ure) & panne, échec — \eng{A power failure stopped the lesson.} « Une coupure de courant a interrompu la leçon. » \\
\hline
failing & adjectif & défaillant — \eng{The failing hard drive must be replaced.} « Le disque dur défaillant doit être remplacé. » \\
\hline
to break down & verbe à particule (broke down, broken down) & tomber en panne — \eng{The server broke down last night.} « Le serveur est tombé en panne hier soir. » \\
\hline
breakdown & nom composé & panne — \eng{The breakdown lasted two hours.} « La panne a duré deux heures. » \\
\hline
broken & adjectif (participe passé) & cassé, hors service — \eng{This keyboard is broken.} « Ce clavier est cassé. » \\
\hline
\end{tabularx}
\end{center}

\begin{attention}[\eng{Break down} : verbe ou nom ?]
Le verbe à particule s'écrit en \textbf{deux mots} : \eng{The computer \textbf{broke down}.} L'ordinateur est tombé en panne. Le nom composé s'écrit en \textbf{un seul mot} : \eng{We had a serious \textbf{breakdown}.} Nous avons eu une panne sérieuse. Ne confondez jamais les deux : \eng{*a break down} (nom) et \eng{*it breakdown} (verbe) sont fautifs.
\end{attention}

\textbf{3. La famille de \eng{use}}

\begin{center}
\begin{tabularx}{\linewidth}{|l|l|X|}
\hline
\rowcolor{popblueL} Mot & Nature et formation & Traduction et exemple \\
\hline
to use & verbe & utiliser — \eng{Students use tablets in class.} « Les élèves utilisent des tablettes en classe. » \\
\hline
user & nom (use + -er) & utilisateur — \eng{Every user needs a password.} « Chaque utilisateur a besoin d'un mot de passe. » \\
\hline
useful & adjectif (use + -ful) & utile — \eng{This tutorial is very useful.} « Ce tutoriel est très utile. » \\
\hline
useless & adjectif (use + -less) & inutile — \eng{Without electricity, the computer is useless.} « Sans électricité, l'ordinateur est inutile. » \\
\hline
user-friendly & adjectif composé & convivial, facile d'usage — \eng{This software is very user-friendly.} « Ce logiciel est très convivial. » \\
\hline
usage & nom & usage, utilisation — \eng{The usage of mobile phones is forbidden during exams.} « L'usage des téléphones portables est interdit pendant les examens. » \\
\hline
\end{tabularx}
\end{center}

\begin{propriete}[Les suffixes \eng{-ful} et \eng{-less}]
Le suffixe \cle{-ful} signifie « plein de » : \eng{useful} (plein d'utilité = utile), \eng{helpful} (serviable). Le suffixe \cle{-less} signifie « sans » : \eng{useless} (sans utilité = inutile), \eng{wireless} (sans fil), \eng{helpless} (démuni). Attention à l'orthographe : \eng{-ful} ne prend qu'\textbf{un seul l} en fin de mot (\eng{useful}, pas \eng{*usefull}).
\end{propriete}

\begin{exemplebox}[Exemples traduits]
\eng{The connection failed, so I had to reconnect.} « La connexion a échoué, j'ai donc dû me reconnecter. »

\eng{This software is very user-friendly.} « Ce logiciel est très convivial / facile d'usage. »

\eng{A power failure caused a total breakdown of the network.} « Une coupure de courant a provoqué une panne totale du réseau. »
\end{exemplebox}

\courssub{Les préfixes utiles : \eng{re-}, \eng{dis-}, \eng{un-}, \eng{mal-}}

\begin{propriete}[Le préfixe \eng{re-} : la répétition]
Le préfixe \cle{re-} placé devant un verbe exprime la \textbf{répétition} de l'action : \eng{restart} (redémarrer), \eng{reboot} (réamorcer, redémarrer), \eng{reset} (réinitialiser), \eng{reconnect} (reconnecter), \eng{reinstall} (réinstaller), \eng{rewrite} (réécrire). Il s'écrit généralement \textbf{sans trait d'union} : \eng{reconnect}, pas \eng{*re-connect}.
\end{propriete}

\begin{propriete}[Les préfixes négatifs \eng{dis-} et \eng{un-}, et le préfixe \eng{mal-}]
\begin{itemize}
\item \cle{dis-} exprime l'\textbf{inverse} de l'action (surtout devant les verbes) : \eng{disconnect} (déconnecter), \eng{disable} (désactiver), \eng{disappear} (disparaître).
\item \cle{un-} exprime la \textbf{négation} ou l'action contraire : \eng{uninstall} (désinstaller), \eng{unavailable} (indisponible), \eng{unable} (incapable), \eng{unlock} (déverrouiller).
\item \cle{mal-} exprime le \textbf{mauvais fonctionnement} : \eng{malfunction} (dysfonctionnement), \eng{malware} (logiciel malveillant).
\end{itemize}
\end{propriete}

\begin{exemplebox}[Les préfixes en contexte]
\eng{The file disappeared, so I had to reinstall the program.} « Le fichier a disparu, j'ai donc dû réinstaller le programme. »

\eng{The website is unavailable; please reboot your router and reconnect.} « Le site est indisponible ; veuillez redémarrer votre routeur et vous reconnecter. »

\eng{The technician disabled the malfunctioning application.} « Le technicien a désactivé l'application défectueuse. »
\end{exemplebox}

\begin{savaistu}[D'où vient le mot \eng{bug} ?]
Le mot \eng{bug} signifie littéralement « insecte ». Selon une anecdote célèbre de l'histoire de l'informatique, en 1947, des ingénieurs américains découvrirent qu'une panne d'un des premiers ordinateurs était causée par un véritable insecte coincé dans la machine. Depuis, le mot désigne toute erreur ou tout défaut dans un programme, et le verbe \eng{to debug} signifie « déboguer », c'est-à-dire corriger les erreurs d'un logiciel.
\end{savaistu}

\courssub{Prononciation : dire les mots des TIC comme un anglophone}

\begin{methode}[La méthode des trois temps]
Pour fixer la prononciation d'un mot nouveau, procédez toujours en trois temps :
\begin{enumerate}
\item \textbf{Écouter} le professeur (ou l'enregistrement) prononcer le mot, sans parler ;
\item \textbf{Répéter en chœur} avec la classe, puis individuellement, en veillant à l'accent tonique ;
\item \textbf{Employer} le mot dans une phrase personnelle : \eng{My screen froze during the test.}
\end{enumerate}
Un mot que l'on n'a jamais dit à voix haute reste un mot étranger.
\end{methode}

\begin{center}
\begin{tabularx}{\linewidth}{|l|l|X|}
\hline
\rowcolor{popblueL} Mot & Prononciation (en lettres françaises) & Piège à éviter \\
\hline
bug & « beug » & ne pas dire « bougue » ; le \eng{u} est bref \\
\hline
crash & « krach » & un seul son « ch » final, comme dans « cache » \\
\hline
freeze & « frize » & \eng{ee} se prononce « i » long ; ne pas dire « frèze » \\
\hline
password & « passe-ouerd » & le \eng{w} se prononce « ou » ; accent sur la première syllabe \\
\hline
troubleshoot & « treubol-choute » & accent sur la première syllabe \\
\hline
technician & « tek-ni-channe » & \eng{-cian} se prononce « channe », jamais « si-anne » \\
\hline
user & « iou-zeur » & \eng{u} initial se prononce « iou », pas « ou » \\
\hline
usage & « iou-sidje » & le \eng{s} se prononce « s », pas « z » \\
\hline
\end{tabularx}
\end{center}

\begin{attention}[L'accent tonique change tout]
En anglais, l'accent tonique est essentiel : on dit \eng{PASSword} et \eng{TECHnician}, avec l'accent sur la première syllabe. Contrairement au français, qui accentue la fin des groupes de mots, l'anglais frappe fortement une syllabe et affaiblit les autres. Écoutez : \eng{the CONnection FAILED}. Répétez en tapant du pied sur la syllabe accentuée.
\end{attention}

\courssub{Orthographe : doublement de consonne et mots composés}

\begin{propriete}[Le doublement de la consonne finale]
Un verbe court d'\textbf{une syllabe} terminé par \textbf{une seule voyelle + une seule consonne} double sa consonne finale devant \eng{-ed} et \eng{-ing} :
\begin{itemize}
\item \eng{bug} → \eng{bugged}, \eng{bugging} (le programme bogue) ;
\item \eng{stop} → \eng{stopped} ; \eng{drop} → \eng{dropped} ; \eng{log} → \eng{logged} (\eng{logged in} = connecté à un compte).
\end{itemize}
Pour les verbes de \textbf{deux syllabes}, on ne double que si l'accent porte sur la dernière syllabe : \eng{install} → \eng{installed} (l'accent est sur « -stall »), mais \eng{open} → \eng{opened} (l'accent est sur « o- », donc pas de doublement).
\end{propriete}

\begin{propriete}[Les mots composés et le trait d'union]
\begin{itemize}
\item \textbf{Adjectifs composés} devant un nom : avec trait d'union — \eng{a user-friendly interface} (une interface conviviale), \eng{an up-to-date system} (un système à jour).
\item \textbf{Noms composés} : souvent en un seul mot ou avec trait d'union — \eng{password} (un seul mot), \eng{log-in} (nom, avec trait d'union), \eng{breakdown} (un seul mot).
\item \textbf{Verbes à particule} : toujours en deux mots, sans trait d'union — \eng{to log in}, \eng{to break down}, \eng{to shut down}.
\end{itemize}
Comparez : \eng{Type your \textbf{log-in}} (nom) et \eng{\textbf{Log in} to your account} (verbe).
\end{propriete}

\begin{attention}[\eng{Connection} et non \eng{*connexion}]
Le nom \eng{connection} s'écrit avec \textbf{-ction} en anglais britannique comme en anglais américain moderne. La forme \eng{connexion} avec \eng{-xion} est vieillie et quasiment abandonnée. Ne recopiez jamais l'orthographe française : \eng{*I have no connexion} est une erreur ; écrivez \eng{I have no connection.} « Je n'ai pas de connexion. »
\end{attention}

\courssub{Entraînement guidé}

\begin{exoresolu}[Complétez avec la bonne forme du mot entre parenthèses]
Énoncé — Complétez chaque phrase avec la forme correcte du mot entre parenthèses.

1. The \dots\dots\dots{} failed during my online interview. (connect)
2. This old printer is completely \dots\dots\dots{}; we must buy a new one. (break)
3. The program kept \dots\dots\dots{}, so the technician uninstalled it. (bug, au moment de l'action)
4. After the power cut, I had to \dots\dots\dots{} the router. (connect, avec l'idée de répétition)
5. This application is not \dots\dots\dots{} at all; nobody understands it. (use, adjectif composé)

\tcblower
Solution détaillée :

1. \eng{The \textbf{connection} failed during my online interview.} Il faut un nom (sujet du verbe) : racine \eng{connect} + suffixe \eng{-ion}. Orthographe : \eng{-ction}, pas \eng{*-xion}.

2. \eng{This old printer is completely \textbf{broken}.} Il faut un adjectif : le participe passé irrégulier de \eng{break} (\eng{break – broke – broken}).

3. \eng{The program kept \textbf{bugging}.} Après \eng{keep}, on emploie la forme en \eng{-ing} ; verbe d'une syllabe « voyelle + consonne » : on double le \eng{g} → \eng{bugging}.

4. \eng{I had to \textbf{reconnect} the router.} L'idée de répétition exige le préfixe \eng{re-}, sans trait d'union.

5. \eng{This application is not \textbf{user-friendly} at all.} Adjectif composé avec trait d'union : \eng{user} + \eng{friendly}.
\end{exoresolu}

\begin{exercice}[À vous de jouer]
1. Formez le nom correspondant à chaque verbe : \eng{connect}, \eng{fail}, \eng{break down}, \eng{use}.
2. Ajoutez le bon préfixe (\eng{re-}, \eng{dis-}, \eng{un-}) : \eng{...install}, \eng{...connect}, \eng{...available}, \eng{...start}, \eng{...able}.
3. Écrivez la forme en \eng{-ed} : \eng{bug}, \eng{install}, \eng{stop}, \eng{open}.
4. Traduisez : « Le serveur est tombé en panne, alors le technicien l'a redémarré et a réinitialisé le système. »
5. Prononcez à voix haute, selon la méthode des trois temps : \eng{bug}, \eng{crash}, \eng{freeze}, \eng{password}, \eng{troubleshoot}, \eng{technician}, puis construisez une phrase personnelle avec chacun.
\end{exercice}

\begin{corrige}[Exercice « À vous de jouer »]
1. \eng{connect → connection} ; \eng{fail → failure} ; \eng{break down → breakdown} ; \eng{use → usage} (ou \eng{user} pour la personne).
2. \eng{\textbf{un}install} ; \eng{\textbf{dis}connect} (ou \eng{\textbf{re}connect}) ; \eng{\textbf{un}available} ; \eng{\textbf{re}start} ; \eng{\textbf{un}able} (ou \eng{\textbf{dis}able}).
3. \eng{bug → bugged} (doublement : une syllabe, voyelle + consonne) ; \eng{install → installed} (accent final) ; \eng{stop → stopped} ; \eng{open → opened} (pas de doublement : accent sur la première syllabe).
4. \eng{The server broke down, so the technician restarted it and reset the system.}
5. Prononciations attendues : « beug », « krach », « frize », « passe-ouerd », « treubol-choute », « tek-ni-channe ». Exemple de phrase : \eng{The technician found a bug and fixed the crash.}
\end{corrige}

\begin{aretenir}[L'essentiel de la section]
\begin{itemize}
\item Les familles de mots se construisent avec des suffixes (\eng{connect → connection → connectivity}) et des préfixes (\eng{re-} = répétition, \eng{dis-} et \eng{un-} = inverse ou négation, \eng{mal-} = mauvais fonctionnement).
\item Verbe en deux mots, nom en un seul : \eng{to break down} mais \eng{a breakdown} ; \eng{to log in} mais \eng{a log-in}.
\item Doublement de la consonne : \eng{bug → bugged}, \eng{install → installed}.
\item Orthographe moderne : \eng{connection} avec \eng{-ction}, jamais \eng{*connexion}.
\item Prononciation à retenir : « beug », « krach », « frize », « passe-ouerd », « treubol-choute », « tek-ni-channe » — et méthode des trois temps : écouter, répéter, employer.
\end{itemize}
\end{aretenir}

\coursec{Réemploi en contexte et production guidée}

Dans la section précédente, nous avons étudié les \cle{familles de mots}, la \cle{prononciation} et l'\cle{orthographe} du lexique des difficultés liées aux TIC (\eng{freeze, crash, log in, connection, virus, troubleshoot}...). Connaître un mot, c'est bien ; savoir l'\textbf{employer spontanément} dans une situation réelle, c'est le véritable objectif. Cette dernière section est donc entièrement consacrée à la \textbf{pratique} : dialoguer au téléphone avec un technicien, rédiger un e-mail de réclamation, jouer une scène de service après-vente, raconter un incident personnel et corriger collectivement les erreurs typiques des francophones. C'est ici que le vocabulaire devient le vôtre.

\courssub{Dialogue modèle : appeler le service technique (helpline)}

\begin{textebox}[At the technical support helpline — dialogue modèle]
\textbf{A:} Good morning, technical support. How can I help you?\\
\textbf{B:} Hello, my laptop keeps freezing and I can't log in to my account.\\
\textbf{A:} Have you tried to restart it?\\
\textbf{B:} Yes, twice, but it crashed again.\\
\textbf{A:} It may be a virus. Run your antivirus and call us back if the problem persists.\\
\textbf{B:} Thank you very much for your help.
\end{textebox}

\textbf{Observons.} Ce court dialogue suit une \textbf{structure fixe} que l'on retrouve dans presque tous les appels à un service d'assistance : salutation et offre d'aide, exposition du problème, questions du technicien, proposition de solution, remerciements. Chaque étape utilise des \textbf{formules ritualisées} qu'il faut mémoriser telles quelles :

\begin{center}
\begin{tabularx}{\linewidth}{|X|X|}
\hline
\rowcolor{popblueL} \textbf{Étape} & \textbf{Formules utiles} \\
\hline
Salutation / offre d'aide & \eng{Good morning, technical support. How can I help you?} \\
\hline
Exposer le problème & \eng{My laptop keeps freezing. / I can't log in. / My screen has gone blank.} \\
\hline
Questions du technicien & \eng{Have you tried to restart it? / When did the problem start?} \\
\hline
Proposer une solution & \eng{Run your antivirus. / Check your connection. / Call us back if the problem persists.} \\
\hline
Remercier et conclure & \eng{Thank you very much for your help. / You're welcome. Goodbye.} \\
\hline
\end{tabularx}
\end{center}

\begin{attention}[Keep + V-ing : une structure à maîtriser]
\eng{My laptop \textbf{keeps freezing}} signifie « mon ordinateur portable \textbf{n'arrête pas de} se bloquer ». La structure est \textbf{keep + verbe en -ing} pour exprimer une action qui se répète de façon agaçante : \eng{The connection keeps dropping.} « La connexion n'arrête pas de tomber. » Ne dites jamais \eng{keeps to freeze}.
\end{attention}

\begin{popfigure}
\begin{tikzpicture}[
  node distance=0.55cm,
  stepS/.style={rectangle, rounded corners, draw=popblue, fill=popblueL, text width=4.6cm, align=center, font=\small, inner sep=5pt},
  arr/.style={-{Stealth[length=2.5mm]}, thick, popdark}
]
\node[stepS, align=center] (s1){\textbf{1. Greeting}\\ \eng{How can I help you?}};
\node[stepS, below=of s1, align=center] (s2){\textbf{2. State the problem}\\ \eng{My laptop keeps freezing.}};
\node[stepS, below=of s2, align=center] (s3){\textbf{3. Describe symptoms}\\ \eng{I can't log in; it crashed twice.}};
\node[stepS, below=of s3, align=center] (s4){\textbf{4. Technician's questions}\\ \eng{Have you tried to restart it?}};
\node[stepS, below=of s4, align=center] (s5){\textbf{5. Solution}\\ \eng{Run your antivirus.}};
\node[stepS, below=of s5, draw=popgreen, fill=popgreenL, align=center] (s6){\textbf{6. Thanks}\\ \eng{Thank you for your help.}};
\draw[arr] (s1) -- (s2);
\draw[arr] (s2) -- (s3);
\draw[arr] (s3) -- (s4);
\draw[arr] (s4) -- (s5);
\draw[arr] (s5) -- (s6);
\end{tikzpicture}
\legende{Organigramme du dialogue d'appel au support technique : six étapes à suivre dans l'ordre.}
\end{popfigure}

\courssub{Texte à trous : un e-mail de réclamation à un cybercafé}

Savoir \textbf{écrire} une réclamation polie mais ferme est une compétence très utile (et souvent demandée à l'examen). Voici un e-mail à compléter avec le lexique étudié.

\begin{exoresolu}[Fill in the gaps — a complaint e-mail]
Énoncé. Complétez cet e-mail avec les mots suivants : \eng{connection — crashed — refund — log in — assistance — disappointed}.

\emph{Dear Sir, I am writing to complain about the service at your cybercafé yesterday. First, the Internet (1) \dots\ was extremely slow, and it was impossible to (2) \dots\ to my e-mail account. Then the computer I was using (3) \dots\ in the middle of my work and I lost my document. I asked for (4) \dots\ , but nobody helped me. I am very (5) \dots\ with your service and I would like a (6) \dots\ . Yours faithfully, A. Mbarga.}

\tcblower
Solution détaillée.
(1) \eng{connection} — la \textbf{connexion} Internet était lente (nom, précédé de l'adjectif \eng{slow}).
(2) \eng{log in} — se \textbf{connecter} à un compte (structure \eng{it was impossible to + base verbale}).
(3) \eng{crashed} — l'ordinateur a \textbf{planté} ; passé simple régulier car le récit est au passé (\eng{yesterday}).
(4) \eng{assistance} — demander de l'\textbf{aide} (nom incountable, jamais \eng{an assistance}).
(5) \eng{disappointed} — je suis \textbf{déçu} (adjectif en \emph{-ed} pour un sentiment ressenti).
(6) \eng{refund} — un \textbf{remboursement} (nom comptable précédé de \eng{a}).

E-mail complet : \eng{Dear Sir, I am writing to complain about the service at your cybercafé yesterday. First, the Internet connection was extremely slow, and it was impossible to log in to my e-mail account. Then the computer I was using crashed in the middle of my work and I lost my document. I asked for assistance, but nobody helped me. I am very disappointed with your service and I would like a refund. Yours faithfully, A. Mbarga.}

« Monsieur, je vous écris pour me plaindre du service de votre cybercafé hier. D'abord, la connexion Internet était extrêmement lente et il était impossible de me connecter à ma messagerie. Ensuite, l'ordinateur que j'utilisais a planté au milieu de mon travail et j'ai perdu mon document. J'ai demandé de l'aide, mais personne ne m'a aidé. Je suis très déçu de votre service et je voudrais un remboursement. Cordialement, A. Mbarga. »
\end{exoresolu}

\begin{attention}[Faux amis et calques dans la réclamation]
\begin{itemize}
\item \eng{assistance} = aide (correct), mais on dit \eng{ask for \textbf{help}} dans la langue courante ; \eng{assistance} est un terme plus formel / soutenu.
\item « Demander un remboursement » = \eng{ask for a refund}, jamais \eng{a reimbursement of my money back} (pléonasme).
\item « Je suis déçu » = \eng{I am disappointed} ; \eng{deceived} signifie « trompé, dupé » (faux ami partiel).
\item Attention à la forme du verbe après \eng{to} : \eng{impossible to \textbf{log in}} (base verbale), pas \eng{to logging in}.
\end{itemize}
\end{attention}

\courssub{Jeu de rôle en binômes : client mécontent / technicien}

\begin{experience}[Role-play — At the repair shop]
Travaillez en binômes. L'élève A est un \textbf{client mécontent} dont l'ordinateur, la tablette ou le téléphone présente un problème (écran noir, batterie qui ne charge plus, connexion impossible, virus, fichiers perdus). L'élève B est le \textbf{technicien} qui pose des questions et propose au moins deux solutions.

\textbf{Étape 1.} Choisissez un problème dans la liste et préparez vos rôles pendant trois minutes : le client note ses symptômes, le technicien note ses questions.

\textbf{Étape 2.} Jouez la scène en suivant l'organigramme à six étapes vu plus haut. Le client doit utiliser au moins trois expressions du lexique (\eng{keeps freezing, crashed, log in, connection dropped, lost my files}) ; le technicien doit proposer des solutions avec l'impératif (\eng{Restart it. Run your antivirus. Check the cable.}).

\textbf{Étape 3.} Inversez les rôles avec un nouveau problème.

\textbf{Étape 4.} Deux binômes volontaires jouent leur scène devant la classe ; les autres élèves relèvent les mots du lexique entendus et signalent les faux amis éventuels.
\end{experience}

\begin{exemplebox}[Exemples de répliques à réemployer]
\eng{My phone keeps switching off by itself.} « Mon téléphone n'arrête pas de s'éteindre tout seul. »

\eng{I lost all my files when the computer crashed.} « J'ai perdu tous mes fichiers quand l'ordinateur a planté. »

\eng{Have you checked the power cable?} « Avez-vous vérifié le câble d'alimentation ? »

\eng{If the problem persists, bring it back and we will repair it free of charge.} « Si le problème persiste, rapportez-le et nous le réparerons gratuitement. »
\end{exemplebox}

\courssub{Rédaction guidée : « A difficulty I faced with ICTs and how I solved it »}

\begin{methode}[Raconter un incident TIC en cinq étapes]
\begin{enumerate}
\item \textbf{Situer} : précisez quand et où (\eng{Last week, during my online lesson... / Two months ago, at the cybercafé in Mokolo...}).
\item \textbf{Décrire le problème} avec le lexique précis (\eng{the connection dropped, my laptop froze, I couldn't log in, the printer was out of order}).
\item \textbf{Expliquer les conséquences} (\eng{I missed half of the lesson. / I lost my document. / I couldn't submit my homework.}).
\item \textbf{Dire la solution trouvée} (\eng{I restarted the router. / A friend helped me. / The technician fixed it.}).
\item \textbf{Conclure} avec un sentiment ou une leçon apprise (\eng{I was relieved. / Now I always save my work. / I learnt to be patient.}).
\end{enumerate}
\end{methode}

\begin{exemplebox}[Modèle de paragraphe (8 à 10 lignes)]
\eng{Last month, I had an online English test at home in Yaoundé. Ten minutes before the test, my laptop suddenly froze and the screen went black. I tried to restart it, but it crashed again. I was very stressed because I thought I would miss the test. Fortunately, my elder brother checked the charger and discovered that the battery was simply flat. He plugged in the laptop and it started normally. I logged in just in time and passed my test. I was so relieved! Since that day, I always check my battery before an online activity.}

« Le mois dernier, j'avais un test d'anglais en ligne à la maison, à Yaoundé. Dix minutes avant le test, mon ordinateur portable s'est soudain bloqué et l'écran est devenu noir. J'ai essayé de le redémarrer, mais il a de nouveau planté. J'étais très stressé car je pensais rater le test. Heureusement, mon grand frère a vérifié le chargeur et a découvert que la batterie était tout simplement vide. Il a branché l'ordinateur et il a démarré normalement. Je me suis connecté juste à temps et j'ai réussi mon test. J'étais tellement soulagé ! Depuis ce jour, je vérifie toujours ma batterie avant une activité en ligne. »
\end{exemplebox}

\begin{exemplebox}[Deux phrases modèles traduites]
\eng{Last week, during my online lesson, the connection dropped three times.} « La semaine dernière, pendant mon cours en ligne, la connexion est tombée trois fois. »

\eng{Fortunately, the technician fixed the problem quickly.} « Heureusement, le technicien a réparé le problème rapidement. »
\end{exemplebox}

\begin{attention}[Au passé, attention aux formes irrégulières]
Le récit d'un incident se fait au \textbf{passé simple}. Mémorisez ces formes irrégulières du lexique des TIC :

\begin{center}
\begin{tabular}{|l|l|l|l|}
\hline
\rowcolor{popblueL} \textbf{Base verbale} & \textbf{Prétérit} & \textbf{Participe passé} & \textbf{Traduction} \\
\hline
freeze & froze & frozen & se bloquer, geler \\
\hline
lose & lost & lost & perdre \\
\hline
forget & forgot & forgotten & oublier \\
\hline
\end{tabular}
\end{center}

Ne dites jamais \eng{my laptop \textbf{freezed}} : la bonne forme est \eng{my laptop \textbf{froze}}. De même : \eng{I \textbf{lost} my files} (pas \eng{losed}) et \eng{I \textbf{forgot} my password} (pas \eng{forgetted}).
\end{attention}

\courssub{Correction collective : chasse aux erreurs}

Pour terminer, la classe corrige ensemble des phrases extraites de productions d'élèves. Chaque phrase contient un faux ami ou une collocation fautive.

\begin{exercice}[Find and correct the mistakes]
Corrigez l'erreur dans chaque phrase, puis justifiez.
\begin{enumerate}
\item \eng{My computer freezed during the examination.}
\item \eng{I could not connect me to my account.}
\item \eng{The connexion of this cybercafé is very slow.}
\item \eng{He gived me a good advise about antivirus software.}
\item \eng{I passed two hours to repair my phone.}
\end{enumerate}
\end{exercice}

\begin{corrige}[Exercice — Find and correct the mistakes]
\begin{enumerate}
\item \eng{My computer \textbf{froze} during the examination.} — \eng{freeze} est irrégulier : \emph{freeze, froze, frozen}.
\item \eng{I could not \textbf{log in to} my account.} — « Se connecter » ne se traduit pas par un verbe réfléchi en anglais ; on utilise la collocation \eng{log in to + compte}.
\item \eng{The \textbf{connection} of this cybercafé is very slow.} — Orthographe anglaise standard : \eng{connection} (avec \emph{-ct-}) ; \eng{connexion} est une variante moins courante. Prononciation : « ké-nèk-cheune ».
\item \eng{He \textbf{gave} me a good \textbf{piece of advice} about antivirus software.} — \eng{give} est irrégulier (\emph{gave}) ; \eng{advice} est incountable : on dit \eng{a piece of advice}, jamais \eng{an advise} (\eng{advise} est le verbe, « conseiller »).
\item \eng{I \textbf{spent} two hours \textbf{repairing} my phone.} — « Passer du temps à faire » = \eng{spend + temps + V-ing} ; « réparer » se dit \eng{repair} ou \eng{fix}, pas de calque de « passer du temps pour ».
\end{enumerate}
\end{corrige}

\begin{aretenir}[L'essentiel de la section]
Pour réemployer le lexique des difficultés avec les TIC : suivez l'organigramme en six étapes pour un appel au support (\emph{greeting, state the problem, describe symptoms, answer questions, accept the solution, thank}) ; structurez un récit d'incident en cinq temps (situer, décrire, conséquences, solution, conclusion) ; au passé, utilisez les formes irrégulières correctes (\emph{froze, lost, forgot, gave}) ; et méfiez-vous des calques du français (\eng{log in to}, \eng{a piece of advice}, \eng{spend time V-ing}). Un vocabulaire bien employé en contexte est un vocabulaire définitivement acquis.
\end{aretenir}

\newpage

\coursec{Activité d'intégration}

\courssub{Objectifs de l'activité}

À la fin de cette activité d'intégration, tu seras capable de :
\begin{itemize}
\item réemployer à l'oral le vocabulaire des difficultés liées aux TIC au travail et dans les études (\eng{slow connection, crash, log in, reset a password, virus, power cut, troubleshoot}...) ;
\item utiliser des collocations et des expressions idiomatiques de la leçon dans un dialogue spontané mais structuré ;
\item poser des questions de diagnostic en anglais correct (\eng{Have you tried...? When did the problem start?}) ;
\item rédiger un court e-mail de suivi (\eng{follow-up email}) résumant un problème technique et sa solution ;
\item éviter les faux amis et les calques du français (\emph{to assist} ne veut pas dire « assister à », \emph{a delay} n'est pas « un délai » dans tous les contextes).
\end{itemize}

\courssub{Situation}

Tu travailles comme stagiaire dans une entreprise de Douala, la \emph{Cameroon Digital Services Ltd}, qui emploie de nombreux jeunes diplômés. Comme dans beaucoup d'entreprises et d'universités camerounaises, les employés et les étudiants dépendent chaque jour des TIC : Wi-Fi, ordinateurs portables, visioconférences, plateformes d'examen en ligne. Mais les difficultés sont fréquentes : coupures d'électricité, connexions lentes, mots de passe oubliés, virus...

L'entreprise a mis en place un \eng{help desk} (service d'assistance informatique). Aujourd'hui, c'est une journée chargée : plusieurs usagers appellent ou se présentent au bureau du technicien avec des problèmes urgents. Voici l'affiche officielle du service, affichée à l'entrée du bureau :

\begin{textebox}[Cameroon Digital Services Ltd — IT Help Desk Notice]
\textbf{IT HELP DESK — WE ARE HERE TO HELP!}\par
Are you having trouble with your computer, your connection or your account? Before you call, please check the following:\par
1. Is your device plugged in and switched on?\par
2. Have you tried restarting your computer?\par
3. Is the Wi-Fi router working? Check the lights.\par
4. Have you forgotten your password? Use the “Forgot password” link to reset it.\par
If the problem persists, come to Room 12 or call extension 404. Describe your problem clearly: tell us what you were doing, what happened, and what error message you saw.\par
\emph{Opening hours: Monday to Friday, 8 a.m. to 5 p.m.}
\end{textebox}

\begin{definition}[Glossaire de l'affiche]
\begin{itemize}
\item \eng{help desk} (n.) : service d'assistance (informatique) — prononciation « help desk ».
\item \eng{to plug in} (v. particule) : brancher (un appareil).
\item \eng{to switch on} (v. particule) : allumer ; contraire : \eng{to switch off}, éteindre.
\item \eng{router} (n.) : routeur, boîtier qui distribue la connexion Wi-Fi — prononciation « raou-teur ».
\item \eng{to reset a password} (collocation) : réinitialiser un mot de passe.
\item \eng{error message} (n.) : message d'erreur.
\item \eng{extension} (n.) : poste (téléphonique interne). Attention, faux ami : ce n'est pas une « extension » au sens d'agrandissement ici.
\end{itemize}
\end{definition}

\courssub{Consignes}

Travail en binômes. L'élève A joue l'usager (employé ou étudiant), l'élève B joue le technicien du help desk.

\textbf{Tâche 1 — Tirage et préparation (10 min).} Tirez au sort une carte-problème parmi les quatre suivantes :
\begin{itemize}
\item \textbf{Card 1:} \eng{The Wi-Fi is extremely slow during your online exam at the university. You are afraid of running out of time.}
\item \textbf{Card 2:} \eng{You have forgotten your password and you must send an urgent report to your boss before noon.}
\item \textbf{Card 3:} \eng{There is a virus on the office computer. Pop-up windows keep appearing and the machine is very slow.}
\item \textbf{Card 4:} \eng{There was a power cut during an important video conference with a partner in Lagos. You were disconnected twice.}
\end{itemize}
Préparez votre dialogue : l'usager décrit le problème avec précision ; le technicien pose au moins trois questions de diagnostic, puis propose au moins deux solutions.

\textbf{Tâche 2 — Jeu de rôle (2 à 3 min par binôme).} Jouez la scène devant la classe. L'usager commence par saluer et expliquer son problème. Le technicien écoute, pose des questions, reformule le problème, puis propose des solutions. L'usager réagit et remercie. Utilisez au moins \textbf{huit mots ou expressions de la leçon} et soulignez-les mentalement : la classe les comptera.

\textbf{Tâche 3 — E-mail de suivi (15 min, à l'écrit).} Après la scène, rédigez ensemble un \eng{follow-up email} de 80 à 120 mots : l'usager écrit au help desk pour résumer le problème, la solution appliquée et le résultat. Respectez la structure d'un e-mail formel : objet (\eng{Subject}), salutation, corps en deux ou trois paragraphes, formule de politesse finale.

\textbf{Tâche 4 — Vote et évaluation (5 min).} La classe vote pour le dialogue le plus clair et le plus riche en vocabulaire de la leçon. Grille d'évaluation : exactitude du lexique (4 pts), collocations correctes (4 pts), questions de diagnostic bien formées (4 pts), prononciation et fluidité (4 pts), e-mail de suivi (4 pts).

\courssub{Ressources à mobiliser}

\begin{aretenir}[Boîte à outils linguistique]
\textbf{Décrire le problème :} \eng{My computer keeps crashing. / The connection is very slow. / I can't log in to my account. / The screen has frozen. / I was disconnected.}

\textbf{Questions de diagnostic :} \eng{When did the problem start? / What were you doing when it happened? / Have you tried restarting it? / Did you see any error message? / Is it plugged in?}

\textbf{Proposer des solutions :} \eng{You should reset your password. / Why don't you run an antivirus scan? / I suggest you move closer to the router. / Let me check the settings.}

\textbf{Expressions idiomatiques :} \eng{to be down} (être en panne), \eng{to sort out a problem} (régler un problème), \eng{to be back online} (être de nouveau connecté), \eng{it works like a charm} (ça marche à merveille).

\textbf{E-mail formel :} \eng{Dear Mr/Ms..., I am writing to report... / Thank you for your assistance. / I look forward to hearing from you. / Yours sincerely,}
\end{aretenir}

\begin{attention}[Pièges à éviter]
\begin{itemize}
\item Ne dis pas \eng{*I have a problem with my computer since two hours} : dis \eng{I have had this problem for two hours} (present perfect + \eng{for}).
\item \eng{Actually} ne signifie pas « actuellement » mais « en fait » ; pour « actuellement », dis \eng{currently} ou \eng{at the moment}.
\item \eng{To assist} veut dire « aider » ; « assister à une réunion » se dit \eng{to attend a meeting}.
\item Ne traduis pas « coupure de courant » par \eng{*electricity cut} : dis \eng{power cut} ou \eng{power outage}.
\end{itemize}
\end{attention}

\courssub{Proposition de production et corrigé commenté}

\begin{exoresolu}[Modèle de dialogue et d'e-mail (Card 3 : virus on the office computer)]
\textbf{Dialogue modèle :}\par
\textbf{A:} Good morning. I need your help. There is a virus on the office computer. Pop-up windows keep appearing and the machine is terribly slow.\par
\textbf{B:} Good morning. Don't worry, we will sort it out. When did the problem start?\par
\textbf{A:} It started yesterday, just after I downloaded a file from an unknown website.\par
\textbf{B:} I see. Have you opened any suspicious e-mail attachments recently?\par
\textbf{A:} Actually, yes. I clicked on a link in an e-mail this morning.\par
\textbf{B:} That is probably the cause. First, you should disconnect the computer from the network so that the virus does not spread. Then I will run a full antivirus scan and update the security software.\par
\textbf{A:} Will I lose my documents?\par
\textbf{B:} No, but I will back up your files first, just to be safe. Next time, avoid downloading files from unknown websites.\par
\textbf{A:} Thank you very much for your assistance.\par
\textbf{B:} You're welcome. The computer should be back online this afternoon.
\tcblower
\textbf{Commentaire des choix de langue :} le dialogue réemploie les collocations de la leçon (\eng{pop-up windows keep appearing}, \eng{run an antivirus scan}, \eng{back up files}, \eng{be back online}) et l'idiome \eng{to sort it out}. Les questions de diagnostic sont au present perfect (\eng{Have you opened...?}) ou au past simple (\eng{When did the problem start?}), distinction essentielle. Notez l'emploi correct de \eng{actually} (= « en fait ») et de \eng{assistance} (= « aide »), deux faux amis évités.\par
\textbf{Follow-up email modèle :}\par
\eng{Subject: Virus on office computer — problem solved}\par
\eng{Dear Mr Etame,}\par
\eng{I am writing to confirm that the virus on the office computer has been removed. The problem started yesterday after I downloaded a file from an unknown website. Pop-up windows kept appearing and the machine was very slow.}\par
\eng{You disconnected the computer from the network, backed up my files and ran a full antivirus scan. The computer is now back online and it works like a charm.}\par
\eng{Thank you for your quick assistance.}\par
\eng{Yours sincerely,}\par
\eng{Laure Ngo Bell, Marketing Department}\par
\textbf{Commentaire :} l'e-mail respecte la structure formelle (objet, salutation, corps, formule finale), utilise le past simple pour les faits passés et réemploie le lexique de la leçon en contexte écrit.
\end{exoresolu}

\courssub{Bilan de l'activité}

Cette activité t'a permis de mobiliser l'ensemble des compétences de la leçon dans une situation de communication authentique : comprendre une consigne écrite (l'affiche du help desk), décrire un problème technique avec le vocabulaire précis, poser des questions de diagnostic grammaticalement correctes, proposer des solutions avec \eng{should}, \eng{why don't you...} et \eng{I suggest...}, puis passer de l'oral à l'écrit avec un e-mail formel. Le vote final t'a entraîné à évaluer la richesse lexicale et la clarté d'une production orale. Retiens la démarche en trois temps du technicien efficace : \cle{écouter} et questionner, \cle{diagnostiquer}, \cle{proposer une solution} — une compétence utile dans le monde du travail anglophone comme camerounais, où la maîtrise du vocabulaire des TIC est désormais indispensable.

\newpage

\coursec{Méthodes et exercices résolus}

\coursec{Lesson 59 — Vocabulary: Words and expressions related to difficulties with the use of ICTs at work or for learning}

\courssub{Méthode 1 : Identifier et employer le vocabulaire des difficultés liées aux TIC}

\begin{methode}[Comment identifier et employer le vocabulaire des difficultés liées aux TIC]
Pour bien utiliser le lexique des problèmes informatiques (\eng{ICT difficulties}), suis ces étapes :
\begin{enumerate}
\item \textbf{Repère la nature du problème} : est-ce un problème de \cle{matériel} (\eng{hardware} : \eng{a broken screen, a faulty keyboard}), de \cle{logiciel} (\eng{software} : \eng{the program crashed}), de \cle{connexion} (\eng{the network is down, a slow connection}) ou de \cle{compétence de l'utilisateur} (\eng{I don't know how to attach a file}) ?
\item \textbf{Choisis le verbe qui colloque avec le nom} : on dit \eng{to experience/to face difficulties}, \eng{the computer crashed/froze}, \eng{to lose data}, \eng{to delete a file by mistake}, \eng{to be cut off} (au téléphone), \eng{to run out of data/credit} (forfait épuisé).
\item \textbf{Décris la conséquence} avec un connecteur : \eng{so, as a result, therefore, that's why} — \eng{The server was down, so I could not submit my assignment online.}
\item \textbf{Propose une solution} : \eng{You should restart the router. / Why don't you update the software? / Have you tried resetting your password?}
\item \textbf{Adapte le registre} : au travail ou à l'école, reste poli et précis (\eng{I'm afraid the projector is not working. Could someone help me, please?}).
\end{enumerate}
Réflexe : apprends toujours un mot avec sa \cle{collocation} : pas seulement \eng{difficulty}, mais \eng{to have difficulty (in) doing something}.
\end{methode}

\begin{exoresolu}[Compléter un texte avec le vocabulaire des TIC]
\textbf{Énoncé} — Fill in the blanks with the correct words from the box: \emph{crashed, connection, password, deleted, attachment, data, frozen, troubleshoot}.

« Yesterday, during my online English lesson, my internet \ldots\ldots\ldots\ldots suddenly stopped working. Then my computer \ldots\ldots\ldots\ldots and the screen \textbf{was completely} \ldots\ldots\ldots\ldots . I had accidentally \ldots\ldots\ldots\ldots an important file, and I could not open the teacher's \ldots\ldots\ldots\ldots because my phone had run out of \ldots\ldots\ldots\ldots . I also forgot my \ldots\ldots\ldots\ldots , so I could not log in again. My father had to \ldots\ldots\ldots\ldots the whole system. »
\tcblower
\textbf{Solution détaillée}
\begin{itemize}
\item \eng{connection} : on parle d'\eng{internet connection} (collocation obligatoire ; jamais « internet connexion »).
\item \eng{crashed} : un ordinateur \eng{crashes} (plante) ; le verbe régulier se conjugue \eng{crash -- crashed -- crashed}.
\item \eng{frozen} : l'écran « gèle » ; participe passé irrégulier de \eng{freeze -- froze -- frozen}. Ici, la phrase impose \eng{was completely} + adjectif/participe \rightarrow \eng{frozen}.
\item \eng{deleted} : \eng{to delete a file} (supprimer un fichier) ; attention, \eng{delete} et non « erase » pour un fichier numérique.
\item \eng{attachment} : pièce jointe d'un e-mail ; ne pas confondre avec le sentiment d'attachement.
\item \eng{data} : ici, le forfait internet (\eng{mobile data}) ; \eng{to run out of data} = épuiser son forfait.
\item \eng{password} : mot de passe ; un seul mot en anglais, jamais « pass word ».
\item \eng{troubleshoot} : dépanner ; verbe irrégulier \eng{troubleshoot -- troubleshot -- troubleshot}.
\end{itemize}
\textbf{Réponse rédigée} : \eng{connection, crashed, frozen, deleted, attachment, data, password, troubleshoot}.
\textbf{Conclusion} : chaque réponse repose sur une collocation apprise ; au Bac, recopie l'orthographe exacte du mot de la boîte.
\end{exoresolu}

\courssub{Méthode 2 : Réemployer le lexique à l'oral et à l'écrit}

\begin{methode}[Comment réemployer le lexique des TIC dans une production]
Pour réutiliser ce vocabulaire dans un dialogue, une lettre de plainte ou un paragraphe argumentatif :
\begin{enumerate}
\item \textbf{Structure ta production en trois temps} : description du problème (\eng{Last week, the school server went down.}), conséquences (\eng{As a result, we could not access the e-learning platform for two days.}), solution ou demande (\eng{I would be grateful if you could repair it as soon as possible.}).
\item \textbf{Utilise des formules toutes faites} : \eng{I am writing to complain about\ldots}, \eng{The main problem is that\ldots}, \eng{To make matters worse,\ldots}, \eng{I suggest that the school should\ldots}
\item \textbf{Varie le lexique} : ne répète pas \eng{problem} ; alterne avec \eng{issue, difficulty, drawback, breakdown, malfunction}.
\item \textbf{Soigne les connecteurs} : \eng{first of all, moreover, however, consequently, finally}.
\item \textbf{Relis-toi} : vérifie l'accord sujet-verbe (\eng{the computer \emph{doesn't} work}), le -s de la 3\textsuperscript{e} personne et l'orthographe des mots techniques (\eng{website, software, hardware} — invariables, sans -s).
\end{enumerate}
\end{methode}

\begin{exoresolu}[Rédiger un court paragraphe de plainte]
\textbf{Énoncé} — Write a paragraph of about 80 words to the head of your school's computer lab. Explain two ICT difficulties you faced during your online revision and suggest one solution. Use at least four words from the lesson (\emph{server, log in, crash, update, connection\ldots}).
\tcblower
\textbf{Raisonnement} : on suit le plan problème -- conséquence -- solution. On choisit deux difficultés : le serveur en panne et l'impossibilité de se connecter. On emploie le prétérit pour les faits passés et le conditionnel poli pour la demande.
\textbf{Réponse rédigée (modèle)} :
\end{exoresolu}

\begin{exemplebox}[Modèle de lettre de plainte]
\eng{Dear Sir, I am writing to complain about the difficulties we faced in the computer lab last week. First of all, the school server went down twice, so we could not log in to the e-learning platform to revise for our exams. Moreover, three computers crashed because their antivirus software has not been updated for months. As a result, many students lost their work. I suggest that the school should update the software and improve the internet connection as soon as possible. Thank you for your attention.}
\end{exemplebox}
« Monsieur, je vous écris pour me plaindre des difficultés rencontrées dans la salle informatique la semaine dernière. D'abord, le serveur de l'école est tombé en panne deux fois, si bien que nous n'avons pas pu nous connecter à la plateforme d'apprentissage en ligne pour réviser nos examens. De plus, trois ordinateurs ont planté parce que leur antivirus n'a pas été mis à jour depuis des mois. Par conséquent, beaucoup d'élèves ont perdu leur travail. Je suggère que l'école mette à jour les logiciels et améliore la connexion internet dès que possible. Merci de votre attention. »
\textbf{Justifications grammaticales} : \eng{went down} = prétérit irrégulier de \eng{go down} (tomber en panne) ; \eng{could not log in} = prétérit de capacité ; \eng{I suggest that the school \emph{should update}} = subjonctif avec \eng{should} après \eng{suggest}.
\textbf{Conclusion} : environ 85 mots, cinq mots du lexique réemployés, registre formel adapté au destinataire.


\courssub{Méthode 3 : Éviter les faux amis et les calques du français}

\begin{methode}[Comment repérer et éviter les faux amis du domaine informatique]
\begin{enumerate}
\item \textbf{Méfie-toi des mots identiques au français} : \eng{actually} signifie « en fait » (pas « actuellement » = \eng{currently}) ; \eng{to assist} signifie « aider » (« assister à un cours en ligne » = \eng{to attend an online class}).
\item \textbf{Ne traduis pas mot à mot} : « j'ai des difficultés à utiliser » = \eng{I have difficulty \emph{using}} (gérondif en -ing, jamais « to use ») ; « être connecté » = \eng{to be connected/logged in}, pas « connected in ».
\item \textbf{Apprends les paires piégeuses par cœur} : \eng{library} (bibliothèque) contre « librairie » (\eng{bookshop}) ; \eng{to achieve} (accomplir) contre « achever » au sens de finir (\eng{to finish}) ; \eng{data} (données/forfait) contre « date ».
\item \textbf{Vérifie la préposition} : \eng{to depend \emph{on}}, \eng{to listen \emph{to}}, \eng{to wait \emph{for}}, \eng{to log \emph{in/into}}, \eng{to suffer \emph{from}} — jamais la préposition française calquée.
\item \textbf{Teste ta phrase} : si tu l'as construite en pensant en français, reformule avec une structure anglaise simple : sujet + verbe + complément.
\end{enumerate}
\end{methode}

\begin{exoresolu}[Corriger les faux amis et les calques]
\textbf{Énoncé} — Each sentence contains one mistake (faux ami or calque du français). Correct it.
\begin{enumerate}
\item \eng{I actually attend online classes every evening.}
\item \eng{She assisted the webinar about cybersecurity last Monday.}
\item \eng{I have difficulty to download the application.}
\item \eng{We are waiting the technician since two hours.}
\item \eng{He depends of his smartphone for everything.}
\end{enumerate}
\tcblower
\textbf{Solution détaillée}
\begin{enumerate}
\item Faux ami : \eng{actually} = « en fait ». Pour « actuellement », il faut \eng{currently}. Correction : \eng{I \emph{currently} attend online classes every evening.} (« Je suis actuellement des cours en ligne chaque soir. »)
\item Faux ami : \eng{to assist} = « aider ». « Assister à » = \eng{to attend}. Correction : \eng{She \emph{attended} the webinar about cybersecurity last Monday.} (prétérit car \eng{last Monday}).
\item Calque : après \eng{have difficulty (in)}, l'anglais exige le gérondif. Correction : \eng{I have difficulty \emph{downloading} the application.}
\item Double erreur : \eng{to wait} exige \eng{for} ; et avec \eng{since/for + durée} pour une action qui continue, on emploie le present perfect continu. Correction : \eng{We \emph{have been waiting for} the technician \emph{for} two hours.} (« since two hours » est un calque de « depuis deux heures » : une durée prend \eng{for}).
\item Préposition calquée : \eng{to depend on}. Correction : \eng{He depends \emph{on} his smartphone for everything.}
\end{enumerate}
\textbf{Conclusion} : à l'examen, relis chaque phrase en te demandant : « Ai-je traduit mot à mot depuis le français ? » Si oui, reformule.
\end{exoresolu}

\courssub{Méthode 4 : Exploiter les familles de mots, la prononciation et l'orthographe}

\begin{methode}[Comment former et orthographier les mots de la même famille]
\begin{enumerate}
\item \textbf{Identifie la nature du mot demandé} dans la phrase (nom, verbe, adjectif, adverbe) grâce à sa place : après \eng{a/the} il faut un nom ; après le sujet, un verbe ; avant un nom, un adjectif ; pour modifier un verbe, un adverbe.
\item \textbf{Applique les suffixes typiques} : \eng{-tion/-sion} (nom : \eng{connect $\rightarrow$ connection}), \eng{-er/-or} (agent : \eng{compute $\rightarrow$ computer}), \eng{-ful/-less} (adjectifs : \eng{use $\rightarrow$ useful/useless}), \eng{-able} (\eng{access $\rightarrow$ accessible}), \eng{-ity} (\eng{secure $\rightarrow$ security}), \eng{-ly} (adverbe : \eng{slow $\rightarrow$ slowly}).
\item \textbf{Attention aux changements orthographiques} : \eng{apply $\rightarrow$ application} (le -y disparaît), \eng{rely $\rightarrow$ reliable}, \eng{technology $\rightarrow$ technological}.
\item \textbf{Mémorise l'orthographe des mots-clés} : \eng{software, hardware} (invariables), \eng{website} (un seul mot), \eng{online/offline} (sans trait d'union devant un nom : \eng{an online course}), \eng{e-mail/email}, \eng{Wi-Fi}.
\item \textbf{Prononciation} : \eng{technology} se dit « tèk-no-lo-dji », \eng{device} « di-vaïce », \eng{data} « déi-ta », \eng{access} « ak-sèss » (accent sur la première syllabe pour le nom).
\end{enumerate}
\end{methode}

\begin{exoresolu}[Word formation : transformer les mots]
\textbf{Énoncé} — Complete each sentence with the correct form of the word in brackets.
\begin{enumerate}
\item The \ldots\ldots\ldots\ldots of the new network took three days. (install)
\item Online learning is very \ldots\ldots\ldots\ldots for students in rural areas. (use)
\item A slow connection makes videoconferences almost \ldots\ldots\ldots\ldots . (possible -- negative)
\item The \ldots\ldots\ldots\ldots of our personal data must be protected. (secure)
\item The technician spoke very \ldots\ldots\ldots\ldots and solved the problem. (clear)
\end{enumerate}
\tcblower
\textbf{Solution détaillée}
\begin{enumerate}
\item Après \eng{the}, il faut un nom : \eng{install $\rightarrow$ installation} (le verbe \eng{install} s'écrit avec deux \eng{l}, le nom \eng{installation} en conserve deux \eng{l}). \eng{The installation of the new network took three days.} (« L'installation du nouveau réseau a pris trois jours. »)
\item Avant le nom sous-entendu, structure \eng{is very + adjectif} : \eng{use $\rightarrow$ useful}. \eng{Online learning is very useful for students in rural areas.}
\item Le sens exige un adjectif négatif (préfixe \eng{im-} devant p) : \eng{possible $\rightarrow$ impossible}. \eng{A slow connection makes videoconferences almost impossible.}
\item Après \eng{the}, un nom : \eng{secure $\rightarrow$ security}. \eng{The security of our personal data must be protected.}
\item Pour modifier le verbe \eng{spoke}, un adverbe : \eng{clear $\rightarrow$ clearly}. \eng{The technician spoke very clearly and solved the problem.}
\end{enumerate}
\textbf{Conclusion} : la place du mot dans la phrase dicte sa nature ; le suffixe donne la forme ; la relecture garantit l'orthographe.
\end{exoresolu}

\courssub{Méthode 5 : Comprendre un texte ou un dialogue sur les difficultés numériques}

\begin{methode}[Comment réussir les questions de compréhension sur ce thème]
\begin{enumerate}
\item \textbf{Lis d'abord les questions}, puis le texte : tu sais quelles informations chercher (\eng{What problem did the student face? What solution was suggested?}).
\item \textbf{Repère les mots du champ lexical} dans le texte : \eng{breakdown, outage, glitch, bug, to freeze, to crash, slow bandwidth, power cut} — ils signalent le passage utile.
\item \textbf{Réponds avec tes propres mots} en reprenant la structure de la question ; ne recopie jamais un long passage sans le reformuler.
\item \textbf{Pour les questions de vocabulaire} (\eng{Find in the text a word meaning\ldots}), cherche un synonyme dans le paragraphe indiqué et recopie-le sans faute.
\item \textbf{Pour les questions vrai/faux}, justifie toujours par une courte citation entre guillemets anglais “ \ldots ”.
\end{enumerate}
\end{methode}

\begin{exoresolu}[Compréhension de texte type Bac]
\textbf{Énoncé} — Read the text and answer the questions.
\begin{textebox}[Digital learning platform launch in Yaoundé]
“When the Ministry launched the digital learning platform, many students in Yaoundé were excited. However, the first week was a nightmare. The platform crashed several times a day because too many users logged in at the same time. In addition, frequent power cuts interrupted the evening classes, and students in remote areas complained about the weak network signal. “I lost my online test twice,” said Marie, a Terminale student. The engineers promised to increase the server capacity and to create an offline version of the application.”
\end{textebox}
Questions : (a) Give two difficulties mentioned in the text. (b) Find in the text a word or expression meaning « pannes d'électricité ». (c) True or false? Justify: « Marie failed her test because she was ill. » (d) What solutions did the engineers promise?
\tcblower
\textbf{Solution détaillée}
\begin{itemize}
\item[(a)] Le texte énumère trois problèmes ; on en choisit deux et on reformule : \eng{The platform crashed several times a day because too many users logged in at the same time, and frequent power cuts interrupted the evening classes.} (On pouvait aussi citer \eng{the weak network signal in remote areas}.)
\item[(b)] \eng{Power cuts} (paragraphe unique) — « pannes d'électricité ». Recopier sans faute, avec le -s du pluriel.
\item[(c)] \eng{False.} Justification par citation : “I lost my online test twice” — \eng{Marie lost her test because of technical problems, not because she was ill.} Le texte ne mentionne aucune maladie.
\item[(d)] \eng{The engineers promised to increase the server capacity and to create an offline version of the application.} (Deux solutions, reliées par \eng{and} ; on garde l'infinitif après \eng{promised to}.)
\end{itemize}
\textbf{Conclusion} : réponses courtes, exactes, appuyées sur le texte ; c'est la stratégie qui rapporte tous les points en compréhension écrite.
\end{exoresolu}

\begin{attention}[Les 5 erreurs qui coûtent des points au Bac]
\begin{enumerate}
\item \textbf{Faux amis} : écrire \eng{I am actually learning online} pour « j'apprends actuellement en ligne ». Correct : \eng{I am \emph{currently} learning online.} De même, « assister à un cours » = \eng{attend a class}, jamais \eng{assist a class}.
\item \textbf{Calque de la construction française} : \eng{I have difficulty to understand} est faux. Après \eng{have difficulty/trouble (in)}, toujours le gérondif : \eng{I have difficulty \emph{understanding} the instructions.}
\item \textbf{Prépositions calquées} : \eng{depend of}, \eng{wait the bus}, \eng{listen the radio} sont fautifs. Correct : \eng{depend \emph{on}}, \eng{wait \emph{for}}, \eng{listen \emph{to}}, \eng{log \emph{in/into}}.
\item \textbf{Orthographe des mots techniques} : \eng{softwares} et \eng{informations} n'existent pas (noms invariables) ; on écrit \eng{website} en un seul mot, \eng{password} en un seul mot, et \eng{an online course} sans trait d'union.
\item \textbf{Temps verbaux avec les marqueurs de durée} : \eng{I am waiting for the technician since two hours} est une double faute. Avec une durée qui continue : present perfect continu + \eng{for} : \eng{I \emph{have been waiting for} the technician \emph{for} two hours.}
\end{enumerate}
\end{attention}

\newpage

\coursec{Exercices}

\courssub{Je vérifie mes connaissances}

\begin{exercice}[QCM : le mot juste]
Choisis la réponse correcte (a, b ou c) pour compléter chaque phrase.
\begin{enumerate}
\item The internet \ldots\ is very slow today. \hfill a) connexion \quad b) connection \quad c) connecter
\item I cannot \ldots\ to my e-mail account; I have forgotten my password. \hfill a) log in \quad b) enter inside \quad c) connect myself
\item The computer \ldots\ during the online exam, so I lost all my answers. \hfill a) crashed \quad b) crushed \quad c) cracked
\item This app is very \ldots; even my grandmother can use it. \hfill a) user-friendly \quad b) friendly-user \quad c) user-friend
\item We had a power \ldots\ in the middle of the video conference. \hfill a) cut \quad b) cup \quad c) cutlery
\end{enumerate}
\end{exercice}

\begin{exercice}[Vrai ou faux, avec justification]
Lis le texte suivant, puis dis si chaque affirmation est \eng{true} ou \eng{false}. Justifie chaque réponse par une ligne du texte.
\begin{textebox}[ICT problems at a Douala secondary school]
At Lycée de New-Bell, the computer room has twenty machines, but only twelve are working. The others broke down months ago and nobody has repaired them. The internet connection is so slow that students often give up before a page loads. Last term, a virus infected the school server and all the examination results were lost because nobody had backed them up. “We are willing to learn,” says a Form Five student, “but the equipment keeps failing us.”
\end{textebox}
\begin{enumerate}
\item All twenty computers in the school are working.
\item The internet connection is fast and reliable.
\item The examination results were saved on an external drive.
\item The students are not interested in learning with computers.
\end{enumerate}
\end{exercice}

\begin{exercice}[Vocabulaire : problèmes et solutions]
Complète chaque phrase avec le mot ou l'expression qui convient, choisi dans la liste : \eng{bandwidth} — \eng{bug} — \eng{breakdown} — \eng{password} — \eng{screen} — \eng{signal}.
\begin{enumerate}
\item There is a \ldots\ in the new software; it closes every time I open a file.
\item I cannot make a call here: there is no \ldots\ in this part of the village.
\item My \ldots\ froze during the lesson, so I had to restart the computer.
\item You need a strong \ldots\ to protect your account from hackers.
\item Watching videos requires a lot of \ldots; our connection is too weak.
\item After the \ldots\ of the main server, the whole office stopped working.
\end{enumerate}
\end{exercice}

\begin{exercice}[Conjugaison et formes verbales]
Mets le verbe entre parenthèses à la forme qui convient (temps, voix ou forme en \emph{-ing}).
\begin{enumerate}
\item Yesterday, the network \ldots\ (to go) down for three hours.
\item I \ldots\ (never / to manage) to download this file; it is too large.
\item The technician \ldots\ (to repair) our computers at the moment.
\item All the data \ldots\ (to lose) because nobody had made a backup.
\item \ldots\ (to restart) the computer often solves the problem.
\end{enumerate}
\end{exercice}

\courssub{Je m'entraîne}

\begin{exercice}[Traduction]
Traduis les phrases suivantes en anglais, en utilisant le lexique de la leçon.
\begin{enumerate}
\item Je n'arrive pas à me connecter : j'ai oublié mon mot de passe.
\item La connexion internet est tombée en panne pendant le cours en ligne.
\item Ce logiciel n'est pas du tout convivial ; il est trop compliqué.
\item N'oublie pas de sauvegarder tes fichiers avant de fermer l'ordinateur.
\item Il y a eu une coupure de courant au milieu de ma visioconférence.
\end{enumerate}
\end{exercice}

\begin{exercice}[Faux amis : corrige les erreurs]
Dans chaque phrase d'élève, souligne le faux ami ou le calque du français, puis réécris la phrase correctement.
\begin{enumerate}
\item I am actually learning computers at school.
\item I assisted the lesson online yesterday.
\item The connexion is very bad in my quarter.
\item I passed two hours on the computer last night.
\item My laptop is in panne; can you help me?
\end{enumerate}
\end{exercice}

\begin{exercice}[Texte à trous : la fracture numérique]
Complète l'article suivant avec les mots et expressions de la liste : \eng{to back up} — \eng{bandwidth} — \eng{breakdown} — \eng{bug} — \eng{to log in} — \eng{user-friendly}.
\begin{textebox}[The digital divide in Cameroon]
In many schools across Cameroon, students face serious difficulties with ICTs. In rural areas, the available \ldots\ is so limited that a single video can take an hour to load. When there is a \ldots\ of the electricity network, all the computers stop at once. Some learning platforms are not \ldots\ at all, and pupils waste precious minutes trying \ldots\ because the instructions are unclear. A small \ldots\ in the examination software once prevented a whole class from submitting their answers. Teachers now advise students \ldots\ their work regularly on a flash drive so that nothing is lost.
\end{textebox}
\end{exercice}

\begin{exercice}[Appariement : problème et solution]
Relie chaque problème (colonne A) à la solution qui convient (colonne B).
\begin{center}
\begin{tabularx}{\linewidth}{|X|X|}
\hline
\rowcolor{popblueL} \textbf{A — Problem} & \textbf{B — Solution} \\
\hline
1. The screen froze. & a) Move closer to the router. \\
\hline
2. There is no Wi-Fi signal. & b) Restart the computer. \\
\hline
3. I have forgotten my password. & c) Charge the battery. \\
\hline
4. The laptop switched off suddenly. & d) Reset it by e-mail. \\
\hline
5. The file will not open. & e) Install the right program. \\
\hline
\end{tabularx}
\end{center}
\end{exercice}

\courssub{Je me prépare au Bac}

\begin{exercice}[Compréhension et langue — type Bac]
Lis attentivement le texte suivant, puis réponds aux questions.
\begin{textebox}[When technology fails the worker]
Manka is an accountant in a small company in Bamenda. Every month, she must send her reports to the head office in Yaoundé through an online platform. Last month, everything went wrong. First, her computer crashed twice on the day of the deadline, and she lost half of her work because she had not saved it. Then, when she finally managed to log in, the platform displayed an error message she could not understand. She called the IT support line, but after twenty minutes of hold music, the line went dead. In the end, she travelled to a cybercafé in town, paid for a faster connection, and sent the report at midnight, one hour before the deadline. “Technology is supposed to make our lives easier,” she says with a tired smile, “but sometimes it makes them much harder.”
\end{textebox}
\textbf{Comprehension questions}
\begin{enumerate}
\item What is Manka's job and where does she work?
\item Give two technical problems she faced on the day of the deadline.
\item Why did she lose half of her work?
\item How did she finally manage to send her report?
\item Do you agree with Manka's last sentence? Justify your answer in two or three sentences.
\end{enumerate}
\textbf{Language work}
\begin{enumerate}
\item Find in the text: (a) a phrasal verb meaning “to stop functioning”, (b) a word meaning “the date limit”, (c) an expression meaning “to enter an account”.
\item Rewrite in the passive voice: “She lost half of her work.”
\item Put into reported speech: “Technology is supposed to make our lives easier,” she says.
\end{enumerate}
\end{exercice}

\begin{exercice}[Traduction — type Bac]
Traduis le passage suivant en anglais.
\begin{textebox}[Texte à traduire]
Hier, pendant le cours d'informatique, le serveur de l'école est tombé en panne. Les élèves n'ont pas pu se connecter à leurs comptes, et ceux qui n'avaient pas sauvegardé leurs fichiers ont tout perdu. Le professeur a expliqué qu'il fallait toujours faire une copie de sécurité avant de fermer sa session. Depuis ce jour, toute la classe fait très attention.
\end{textebox}
\end{exercice}

\begin{exercice}[Composition — type Bac]
Rédige une composition en anglais de \textbf{80 à 100 mots} sur le sujet suivant :
\begin{center}
\emph{“Describe a difficulty you or someone you know faced with ICTs at school or at work, and explain how it was solved.”}
\end{center}
Consignes :
\begin{itemize}
\item Présente brièvement la personne et le contexte (école, travail, cybercafé).
\item Décris la difficulté technique en utilisant au moins \textbf{quatre mots ou expressions} du lexique de la leçon (par exemple : \eng{to crash}, \eng{to log in}, \eng{breakdown}, \eng{to back up}, \eng{user-friendly}, \eng{bug}).
\item Explique la solution trouvée et conclus par une phrase d'opinion personnelle.
\end{itemize}
\end{exercice}



\newpage

\coursec{Corrigés des exercices}

\begin{corrige}[Exercice 1 — QCM : le mot juste]
\begin{enumerate}
\item \textbf{b) connection}. Orthographe anglaise : \eng{connection} (avec \emph{-ect-}). \eng{Connexion} est un calque du français (forme rare et vieillie en anglais, à éviter).
\item \textbf{a) log in}. Phrasal verb : \eng{to log in to an account} (se connecter à un compte). \eng{Enter inside} est un pléonasme ; \eng{connect myself} est un calque de « me connecter ».
\item \textbf{a) crashed}. \eng{To crash} = (ordinateur) planter, cesser de fonctionner. \eng{To crush} = écraser ; \eng{to crack} = fissurer.
\item \textbf{a) user-friendly}. Adjectif composé d'ordre fixe : « convivial, facile à utiliser ».
\item \textbf{a) cut}. \eng{A power cut} = une coupure de courant. \eng{Cup} = tasse ; \eng{cutlery} = les couverts.
\end{enumerate}
Barème indicatif : 1 point par bonne réponse (5 points).
\end{corrige}

\begin{corrige}[Exercice 2 — Vrai ou faux, avec justification]
\begin{enumerate}
\item \textbf{False.} Justification : \eng{“The computer room has twenty machines, but only twelve are working.”}
\item \textbf{False.} Justification : \eng{“The internet connection is so slow that students often give up before a page loads.”}
\item \textbf{False.} Justification : \eng{“All the examination results were lost because nobody had backed them up.”}
\item \textbf{False.} Justification : \eng{“We are willing to learn, but the equipment keeps failing us.”}
\end{enumerate}
Point de méthode : au Bac, une réponse \eng{true} ou \eng{false} sans justification tirée du texte ne rapporte généralement pas tous les points ; recopie toujours la ligne exacte du texte entre guillemets.
\end{corrige}

\begin{corrige}[Exercice 3 — Vocabulaire : problèmes et solutions]
\begin{enumerate}
\item \textbf{bug} — \eng{There is a bug in the new software; it closes every time I open a file.} (« bogue » : défaut dans un programme.)
\item \textbf{signal} — \eng{There is no signal in this part of the village.} (pas de réseau.)
\item \textbf{screen} — \eng{My screen froze during the lesson, so I had to restart the computer.} (\eng{To freeze} : se bloquer ; prétérit \emph{froze}, participe passé \emph{frozen}.)
\item \textbf{password} — \eng{You need a strong password to protect your account from hackers.} (mot de passe.)
\item \textbf{bandwidth} — \eng{Watching videos requires a lot of bandwidth; our connection is too weak.} (bande passante ; nom indénombrable : pas de pluriel.)
\item \textbf{breakdown} — \eng{After the breakdown of the main server, the whole office stopped working.} (panne.)
\end{enumerate}
\end{corrige}

\begin{corrige}[Exercice 4 — Conjugaison et formes verbales]
\begin{enumerate}
\item \textbf{went} — \eng{Yesterday, the network went down for three hours.} Prétérit simple de \eng{to go} (\emph{go – went – gone}), marqué par \eng{yesterday}.
\item \textbf{have never managed} — \eng{I have never managed to download this file.} Present perfect (expérience jusqu'à aujourd'hui) ; \eng{never} se place entre l'auxiliaire \eng{have} et le participe passé.
\item \textbf{is repairing} — \eng{The technician is repairing our computers at the moment.} Present continu (\emph{be + V-ing}), marqué par \eng{at the moment}.
\item \textbf{were lost} — \eng{All the data were lost because nobody had made a backup.} Voix passive au prétérit (\emph{were} + participe passé \emph{lost}) ; \eng{data} est ici traité comme un pluriel.
\item \textbf{Restarting} — \eng{Restarting the computer often solves the problem.} Gérondif (forme en \emph{-ing}) en fonction de sujet de la phrase.
\end{enumerate}
\end{corrige}

\begin{corrige}[Exercice 5 — Traduction]
\begin{enumerate}
\item \eng{I cannot log in: I have forgotten my password.} (Pas de réflexif : \eng{to log in} suffit, jamais \emph{to connect myself}.)
\item \eng{The internet connection went down during the online lesson.} (\eng{To go down} = tomber en panne ; prétérit \emph{went}.)
\item \eng{This software is not user-friendly at all; it is too complicated.} (« Convivial » pour un logiciel = \eng{user-friendly}, jamais \emph{convivial}.)
\item \eng{Don't forget to back up your files before you switch off the computer.} (\eng{To back up} = sauvegarder ; \eng{to switch off} = éteindre.)
\item \eng{There was a power cut in the middle of my video conference.} (\eng{A power cut} = une coupure de courant.)
\end{enumerate}
Barème indicatif : 2 points par phrase (lexique correct 1 point, grammaire correcte 1 point).
\end{corrige}

\begin{corrige}[Exercice 6 — Faux amis : corrige les erreurs]
\begin{enumerate}
\item \eng{actually} $\rightarrow$ \eng{I am currently learning computer skills at school.} (\eng{Actually} = en fait ; « actuellement » = \eng{currently} / \eng{at present}.)
\item \eng{assisted} $\rightarrow$ \eng{I attended the lesson online yesterday.} (\eng{To assist} = aider ; « assister à » = \eng{to attend}.)
\item \eng{connexion} $\rightarrow$ \eng{The connection is very bad in my neighbourhood.} (Orthographe anglaise \eng{connection} ; « quartier » = \eng{neighbourhood}, pas \emph{quarter} dans ce sens.)
\item \eng{passed} $\rightarrow$ \eng{I spent two hours on the computer last night.} (« Passer du temps » = \eng{to spend time} ; \eng{to pass} = réussir un examen ou dépasser.)
\item \eng{in panne} $\rightarrow$ \eng{My laptop has broken down; can you help me?} (« Être en panne » = \eng{to break down} / \eng{to be out of order}.)
\end{enumerate}
\end{corrige}

\begin{corrige}[Exercice 7 — Texte à trous : la fracture numérique]
Dans l'ordre :
\begin{enumerate}
\item \textbf{bandwidth} — \eng{the available bandwidth is so limited} (nom indénombrable, précédé de l'adjectif \eng{available}).
\item \textbf{breakdown} — \eng{a breakdown of the electricity network} (une panne du réseau électrique).
\item \textbf{user-friendly} — \eng{some learning platforms are not user-friendly at all}.
\item \textbf{to log in} — \eng{pupils waste precious minutes trying to log in} (structure : \eng{to try to do something}).
\item \textbf{bug} — \eng{a small bug in the examination software} (article indéfini \eng{a} devant nom comptable).
\item \textbf{to back up} — \eng{teachers now advise students to back up their work} (structure : \eng{to advise somebody to do something}).
\end{enumerate}
\end{corrige}

\begin{corrige}[Exercice 8 — Appariement : problème et solution]
1 — b) : \eng{The screen froze. $\rightarrow$ Restart the computer.}\par
2 — a) : \eng{There is no Wi-Fi signal. $\rightarrow$ Move closer to the router.}\par
3 — d) : \eng{I have forgotten my password. $\rightarrow$ Reset it by e-mail.}\par
4 — c) : \eng{The laptop switched off suddenly. $\rightarrow$ Charge the battery.}\par
5 — e) : \eng{The file will not open. $\rightarrow$ Install the right program.}
\end{corrige}

\begin{corrige}[Exercice 9 — Compréhension et langue — type Bac]
\textbf{Comprehension questions}
\begin{enumerate}
\item \eng{Manka is an accountant. She works in a small company in Bamenda.}
\item \eng{Her computer crashed twice, and the platform displayed an error message she could not understand.} (On accepte aussi : \eng{the IT support line went dead}.)
\item \eng{She lost half of her work because she had not saved it.}
\item \eng{She travelled to a cybercafé in town, paid for a faster connection and sent the report from there, one hour before the deadline.}
\item Réponse personnelle. Exemple : \eng{I agree with Manka. Technology saves time when it works, but breakdowns, slow connections and power cuts can waste a whole day. That is why we should always back up our work.}
\end{enumerate}
\textbf{Language work}
\begin{enumerate}
\item (a) \eng{to crash} (on accepte aussi \eng{to go down} dans \eng{the line went dead}) ; (b) \eng{deadline} ; (c) \eng{to log in}.
\item \eng{Half of her work was lost (by her).} (Voix passive au prétérit : \emph{was} + participe passé \emph{lost}.)
\item \eng{She says (that) technology is supposed to make their lives easier.} (Avec un verbe introducteur au présent, \eng{says}, les temps ne changent pas ; le possessif \eng{our} devient \eng{their}.)
\end{enumerate}
Barème indicatif : compréhension 5 $\times$ 2 points ; langue 3 + 2 + 2 points. Points de vigilance : répondre par des phrases complètes ; ne pas recopier tout un paragraphe ; respecter la concordance des temps.
\end{corrige}

\begin{corrige}[Exercice 10 — Traduction — type Bac]
Proposition de traduction :
\textbf{Version corrigée}\par
\eng{Yesterday, during the computer science lesson, the school server broke down. The students could not log in to their accounts, and those who had not backed up their files lost everything. The teacher explained that it was always necessary to make a backup before logging out. Since that day, the whole class has been very careful.}

Points de vigilance :
\begin{itemize}
\item \eng{to break down} : prétérit irrégulier \emph{broke down} (jamais \emph{breaked}).
\item \eng{to log in to their accounts} : préposition \eng{to} après \eng{log in}.
\item \eng{to back up} (sauvegarder) et \eng{to log out} (fermer sa session) : particules séparables obligatoires.
\item Discours rapporté : \eng{the teacher explained that it was always necessary} (recul du temps présent vers le prétérit).
\item \eng{Since that day} + present perfect : \eng{the whole class has been very careful}.
\end{itemize}
Barème indicatif : 2 points par phrase (lexique 1 point, morphosyntaxe 1 point).
\end{corrige}

\begin{corrige}[Exercice 11 — Composition — type Bac]
Modèle de production (environ 95 mots) :
\textbf{Model answer}\par
\eng{Last term, my cousin Clarisse, who works in a bank in Douala, faced a serious problem. On the day she had to send an important report, her computer crashed and the office server broke down. She could not log in to her account, and she had forgotten to back up her file. Fortunately, a technician repaired the breakdown quickly and recovered her document. She finally sent the report from a cybercafé with a faster connection. In my opinion, this story shows that technology is useful, but we must always save and back up our work.}

Barème indicatif (sur 10) : respect du sujet et du nombre de mots (2 points) ; emploi correct d'au moins quatre termes du lexique de la leçon (4 points) ; exactitude grammaticale, récit au prétérit (2 points) ; organisation et conclusion d'opinion (\eng{in my opinion}, \eng{I think that}) (2 points).
Points de vigilance : raconter au prétérit simple ; éviter les calques (\emph{connexion}, \emph{in panne}, \emph{passed two hours}) ; une seule idée par phrase ; relire pour vérifier les terminaisons en \emph{-ed}.
\end{corrige}

\newpage

\coursec{Fiche bilan}

\begin{popfigure}
\begin{tikzpicture}[mindmap, grow cyclic, every node/.style={concept, text=white, align=center, font=\small},
  root concept/.append style={concept color=popblue, font=\bfseries\small},
  level 1/.append style={level distance=3.4cm, sibling angle=51.4},
  level 2/.append style={level distance=2.1cm, sibling angle=45, font=\scriptsize}]
\node[root concept]{ICT difficulties at work and for learning}
  child[concept color=poporange]{ node{Technical problems}
    child{ node{break down, crash, freeze} }
    child{ node{bug, virus, error message} }
  }
  child[concept color=popgreen]{ node{Connection issues}
    child{ node{slow / unstable network} }
    child{ node{power cut, blackout} }
  }
  child[concept color=poppurple]{ node{User difficulties}
    child{ node{lack digital skills} }
    child{ node{get lost, feel confused} }
  }
  child[concept color=poppink]{ node{Learning online}
    child{ node{distraction, isolation} }
    child{ node{screen fatigue} }
  }
  child[concept color=popgold]{ node{Solutions}
    child{ node{troubleshoot, restart} }
    child{ node{update, back up} }
  }
  child[concept color=popdark]{ node{Collocations}
    child{ node{log in, sign up} }
    child{ node{run out of data} }
  }
  child[concept color=popmuted]{ node{False friends}
    child{ node{actually $\neq$ actuellement} }
    child{ node{library $\neq$ librairie} }
  };
\end{tikzpicture}
\legende{Carte mentale : le lexique des difficultés liées aux TIC au travail et à l'école}
\end{popfigure}

\begin{bilan}
\textbf{1. Lexique clé (English -- Français)}
\begin{itemize}
  \item \eng{to break down} : tomber en panne ; \eng{a breakdown} : une panne
  \item \eng{to crash / to freeze} : planter / se figer (ordinateur)
  \item \eng{a bug} : un bogue ; \eng{a virus} : un virus ; \eng{an error message} : un message d'erreur
  \item \eng{a slow / unstable connection} : une connexion lente / instable
  \item \eng{a power cut / a blackout} : une coupure de courant
  \item \eng{to run out of data / credit} : être à court de forfait internet / de crédit
  \item \eng{to log in / to log out} : se connecter / se déconnecter
  \item \eng{to sign up} : s'inscrire ; \eng{to update} : mettre à jour ; \eng{to back up} : sauvegarder
  \item \eng{to troubleshoot} : dépanner ; \eng{tech support} : l'assistance technique
  \item \eng{digital skills} : compétences numériques ; \eng{the digital divide} : la fracture numérique
  \item \eng{screen fatigue} : fatigue liée aux écrans ; \eng{distraction} : la distraction
\end{itemize}

\textbf{2. Collocations et expressions idiomatiques}
\begin{itemize}
  \item \eng{to have trouble (with) + noun / -ing} : avoir des difficultés (avec / à) -- \eng{I'm having trouble connecting.} « J'ai du mal à me connecter. »
  \item \eng{to be down} : être hors service -- \eng{The network is down.} « Le réseau est en panne. »
  \item \eng{to be at a loss} : être perdu, désemparé
  \item \eng{to get the hang of something} : prendre le coup de main
  \item \eng{it's not rocket science} : ce n'est pas sorcier
\end{itemize}

\textbf{3. Faux amis et calques à éviter}
\begin{center}
\begin{tabularx}{\linewidth}{|X|X|X|}
\hline
\rowcolor{popblueL} Français visé & English correct & Piège à éviter \\
\hline
actuellement & currently, at present & \eng{actually} = en fait \\
\hline
une librairie & a bookshop & \eng{a library} = une bibliothèque \\
\hline
assister à (un cours) & to attend (a class) & \eng{to assist} = aider \\
\hline
une formation & training & \eng{a formation} = une formation (géologique, militaire) \\
\hline
j'ai des difficultés & I have difficulties / trouble & pas de calque \eng{I have difficults} \\
\hline
\end{tabularx}
\end{center}

\textbf{4. Familles de mots et prononciation}
\begin{itemize}
  \item \eng{connect} (v.) $\rightarrow$ \eng{connection} (n.) $\rightarrow$ \eng{connected / disconnected} (adj.)
  \item \eng{difficult} (adj.) $\rightarrow$ \eng{difficulty} (n., plur. \eng{difficulties})
  \item \eng{technology} (n.) $\rightarrow$ \eng{technological} (adj.) $\rightarrow$ \eng{technologically} (adv.)
  \item Prononciation : \eng{technology} « tek-no-lo-dji », \eng{troubleshoot} « treu-beul-choute », \eng{virus} « vaï-reuss », \eng{data} « déï-ta ».
\end{itemize}

\textbf{5. Méthode express}
\begin{enumerate}
  \item Repérer le thème du texte (panne, connexion, compétences, solutions).
  \item Souligner les verbes à particule (\eng{break down, log in, back up}) et noter leur sens en contexte.
  \item Vérifier chaque mot « transparent » : est-ce un faux ami ?
  \item Réemployer 5 mots nouveaux dans des phrases personnelles liées au Cameroun (cybercafé, cours en ligne, coupure de courant).
\end{enumerate}
\end{bilan}

\begin{aretenir}[Checklist avant le Bac]
\begin{itemize}
  \item[$\square$] Je connais au moins 15 mots du champ lexical des difficultés liées aux TIC.
  \item[$\square$] Je sais employer \eng{to break down}, \eng{to crash}, \eng{to freeze} et \eng{to be down} correctement.
  \item[$\square$] Je maîtrise les verbes à particule : \eng{log in, log out, sign up, back up, run out of}.
  \item[$\square$] Je sais dire « avoir du mal à » avec \eng{to have trouble + -ing}.
  \item[$\square$] Je distingue \eng{actually / currently} et \eng{library / bookshop}.
  \item[$\square$] Je ne traduis jamais « assister à un cours » par \eng{to assist a class} mais par \eng{to attend a class}.
  \item[$\square$] Je connais la famille de \eng{connect} et celle de \eng{difficult}.
  \item[$\square$] Je sais prononcer \eng{technology}, \eng{data} et \eng{virus} à l'anglaise.
  \item[$\square$] Je sais proposer des solutions : \eng{troubleshoot, restart, update, contact tech support}.
  \item[$\square$] Je peux rédiger un court paragraphe sur les difficultés des TIC dans mon école avec 5 mots de la leçon.
  \item[$\square$] Je relie cette leçon au chapitre \eng{Media and Communication} pour l'épreuve d'expression écrite.
\end{itemize}
\end{aretenir}

\findecours

\end{document}
